\documentclass[11pt,reqno]{amsart}

% ================================================================
%  Packages
% ================================================================
\usepackage{amsmath,amssymb,amsthm}
\usepackage{mathtools}
\usepackage{mathrsfs}
\usepackage[shortlabels]{enumitem}
\usepackage{booktabs}
\usepackage{array}
\usepackage{microtype}
\usepackage[dvipsnames]{xcolor}
\usepackage[colorlinks=true,
  linkcolor=blue!60!black,
  citecolor=green!40!black,
  urlcolor=blue!60!black]{hyperref}

% ================================================================
%  Theorem environments
% ================================================================
\newtheorem{theorem}{Theorem}[section]
\newtheorem{proposition}[theorem]{Proposition}
\newtheorem{lemma}[theorem]{Lemma}
\newtheorem{corollary}[theorem]{Corollary}
\newtheorem{conjecture}[theorem]{Conjecture}
\theoremstyle{definition}
\newtheorem{definition}[theorem]{Definition}
\newtheorem{construction}[theorem]{Construction}
\newtheorem{example}[theorem]{Example}
\theoremstyle{remark}
\newtheorem{remark}[theorem]{Remark}

% ================================================================
%  Macros
% ================================================================
\newcommand{\cA}{\mathcal{A}}
\newcommand{\cD}{\mathcal{D}}
\newcommand{\cH}{\mathcal{H}}
\newcommand{\cM}{\mathcal{M}}
\newcommand{\cW}{\mathcal{W}}
\newcommand{\barB}{\overline{B}}
\newcommand{\Ran}{\mathrm{Ran}}
\newcommand{\FM}{\overline{\mathrm{FM}}}
\newcommand{\Conf}{\mathrm{Conf}}
\newcommand{\fg}{\mathfrak{g}}
\newcommand{\Sym}{\mathrm{Sym}}
\newcommand{\Res}{\operatorname{Res}}
\newcommand{\Spec}{\operatorname{Spec}}
\newcommand{\Diff}{\mathrm{Diff}}
\newcommand{\Vir}{\mathrm{Vir}}
\newcommand{\Hom}{\mathrm{Hom}}
\newcommand{\id}{\mathrm{id}}
\renewcommand{\Bbbk}{\mathbf{C}}

\DeclareMathOperator{\gr}{gr}
\DeclareMathOperator{\End}{End}

\numberwithin{equation}{section}

% ================================================================
\begin{document}

\title[Modular Koszul Duality: a Programme]
{Modular Koszul Duality for Chiral Algebras\\
on Algebraic Curves: a Programme}

\author{Raeez Lorgat}
\address{Perimeter Institute for Theoretical Physics,
  Waterloo, ON N2L 2Y5, Canada}
\email{rlorgat@perimeterinstitute.ca}

\date{\today}

\begin{abstract}
We survey the modular Koszul duality programme for chiral
algebras on algebraic curves.  The bar construction of a
chiral algebra produces a factorization coalgebra whose
differential encodes operator product singularities and
whose curvature at genus~$g \ge 1$ is controlled by a
single scalar, the modular characteristic~$\kappa$.
Five theorems describe the resulting structure: bar-cobar
adjunction with Verdier intertwining~(A), bar-cobar
inversion on the Koszul locus~(B), complementarity
decomposition with shifted-symplectic Lagrangian
geometry~(C), the genus expansion
$F_g = \kappa \cdot \lambda_g^{\mathrm{FP}}$ (uniform-weight) governed by
the Hirzebruch $\hat{A}$-class~(D), and polynomial
chiral Hochschild cohomology~(H).  All five are projections
of a single universal Maurer--Cartan element~$\Theta_\cA$
in the modular convolution algebra.  The shadow obstruction
tower, the G/L/C/M depth classification, the seven faces
of the collision residue, the Swiss-cheese realization, and
the standard landscape of worked examples are described.
The programme extends to holographic systems, analytic
sewing, and Calabi--Yau geometry in three dimensions; the
single open conjecture is the CY-A correspondence at
$d = 3$.
The programme spans three volumes totalling 4,542 pages,
with 120K+ computational tests; six new proved results
include the universal $\alpha_g$ identity, the soft
graviton coefficient $a_3 = 2$, and free-field exactness.
\end{abstract}

\maketitle

% ====================================================================
\section{The bar construction on algebraic curves}
\label{sec:bar}
% ====================================================================

The bar construction of an associative algebra is classical:
given an augmented algebra~$A$ over a field, form the tensor
coalgebra $T^c(s^{-1}\bar A) = \bigoplus_{n \ge 1}
(s^{-1}\bar A)^{\otimes n}$ on desuspended generators, equip it
with the differential that encodes the multiplication, and
observe that $d^2 = 0$ is equivalent to associativity.
Koszul duality~\cite{BGS96,Priddy70}, homological algebra,
and deformation theory all descend from this single
construction.

A chiral algebra on a smooth algebraic curve~$X$ encodes
operator product singularities: the product of two sections
$a(z)\,b(w)$ diverges as $z \to w$, and the divergence carries
algebraic data.  The problem is to build the bar construction in
this setting.  The point is replaced by a curve; the
multiplication map is replaced by a chiral bracket with poles
along the diagonal; the tensor coalgebra must account for the
geometry of colliding points on~$X$.

% ====================================================================
\subsection{The propagator}
\label{ssec:propagator}
% ====================================================================

On the configuration space $C_2(X) = X \times X \setminus \Delta$
of two distinct points on a smooth algebraic curve~$X$
over~$\Bbbk$, define
\begin{equation}\label{eq:eta}
\eta_{12}
\;=\;
d\log(z_1 - z_2)
\;=\;
\frac{dz_1 - dz_2}{z_1 - z_2}\,.
\end{equation}
This $1$-form has a simple pole along the diagonal
$\Delta \subset X \times X$ with residue~$1$.  The residue is
intrinsic: it depends on no coordinate, no metric, no level.

Three properties single out~$\eta_{12}$.  A double pole
$dz/(z-w)^2$ transforms under coordinate change and has no
coordinate-independent residue.  A regular $1$-form has
residue zero and extracts nothing from the collision.  The
logarithmic derivative is the unique $1$-form on~$C_2(X)$
with a first-order pole along~$\Delta$ and conformally
invariant residue.

On $C_3(X)$, the three pullbacks
$\eta_{12}, \eta_{23}, \eta_{31}$ satisfy
\begin{equation}\label{eq:arnold}
\eta_{12} \wedge \eta_{23}
\;+\;
\eta_{23} \wedge \eta_{31}
\;+\;
\eta_{31} \wedge \eta_{12}
\;=\; 0\,.
\end{equation}
This is the \emph{Arnold relation}.  It restates the
partial-fractions identity
\[
\frac{1}{(z_1 {-} z_2)(z_2 {-} z_3)}
+
\frac{1}{(z_2 {-} z_3)(z_3 {-} z_1)}
+
\frac{1}{(z_3 {-} z_1)(z_1 {-} z_2)}
= 0\,,
\]
which holds because the left side, viewed as a rational function
of~$z_1$ with $z_2, z_3$ fixed, has simple poles at $z_2$
and~$z_3$ whose residues sum to zero.
Arnold~\cite{Arnold69} proved that~\eqref{eq:arnold}, replicated
across all triples $\{i,j,k\} \subset \{1,\ldots,n\}$, generates
the complete ideal of relations in
$H^*(\Conf_n(\Bbbk), \Bbbk)$.

Under a coordinate change $z \mapsto w(z)$, the propagator
transforms as
$\eta_{ij} \mapsto \eta_{ij}
+ d\log\frac{w(z_i) - w(z_j)}{z_i - z_j}$.
Near the diagonal the correction becomes $d\log w'(z_i)$: the
standard cocycle for~$\Diff(X)$.  The propagator is a BRST ghost
for diffeomorphisms; its failure to be coordinate-invariant is
the failure that the Virasoro algebra measures.


% ====================================================================
\subsection{Chiral algebras}
\label{ssec:chiral}
% ====================================================================

A \emph{chiral algebra}~$\cA$ on~$X$ is a
$\cD_X$-module equipped with a chiral bracket
\[
\mu^{\mathrm{ch}}
\;\colon\;
j_* j^*(\cA \boxtimes \cA)
\;\longrightarrow\;
\Delta_!(\cA)\,,
\]
where $j \colon X \times X \setminus \Delta
\hookrightarrow X \times X$ is the open complement and
$\Delta_! \cA$ is the $!$-pushforward along the diagonal
embedding~\cite{BD04}.  The source consists of sections of
$\cA \boxtimes \cA$ with poles of arbitrary order
along~$\Delta$; the chiral bracket extracts the singular
part and projects it onto the diagonal.  In local coordinates,
the chiral bracket encodes the operator product expansion:
\[
a(z)\,b(w)
\;\sim\;
\sum_{n \ge 0}
\frac{a_{(n)}b(w)}{(z - w)^{n+1}}\,.
\]
The $n$-th products $a_{(n)}b \in \cA$ are the Fourier
modes, extracted by
$a_{(n)}b
= \Res_{z=w}[a(z)\,b(w)\,(z{-}w)^n]$.
The zeroth product $a_{(0)}b$ is the Lie bracket
(first-order pole); $a_{(1)}b$ carries the
symmetric/metric data (second-order pole).  Higher products
$a_{(n)}b$ for $n \ge 2$ are determined by the
$\cD_X$-module structure.

The Borcherds identity
\begin{equation}\label{eq:borcherds}
\sum_{j \ge 0} \binom{m}{j}
(a_{(n+j)}b)_{(m+k-j)}c
\;=\;
\sum_{j \ge 0} (-1)^j \binom{n}{j}
\Bigl(
a_{(m+n-j)}(b_{(k+j)}c)
- (-1)^n b_{(n+k-j)}(a_{(m+j)}c)
\Bigr)
\end{equation}
is the complete axiom system.  At $n = 0$ it reduces to the
Jacobi identity for the Lie bracket~$a_{(0)}$; at $n = 1$,
$m = k = 0$, to the derivation property of~$a_{(0)}$ with
respect to~$a_{(1)}b$.


% ====================================================================
\subsection{The bar complex}
\label{ssec:bar-complex}
% ====================================================================

Place sections $a_1, \ldots, a_n \in \cA$ at distinct points
$z_1, \ldots, z_n \in X$.  The Fulton--MacPherson
compactification $\FM_n(X)$ resolves $\Conf_n(X)$ by
iterated real oriented blowup along all diagonal
strata~\cite{FM94}.  Points of
$\FM_n(X) \setminus \Conf_n(X)$ record the limiting direction
and relative scale of approach as points collide.  The boundary
is a manifold with corners; codimension-one faces $D_S$
are indexed by subsets $S \subseteq \{1,\ldots,n\}$ with
$|S| \ge 2$, recording which points have collided.  Each face
factors:
\begin{equation}\label{eq:fm-boundary}
D_S
\;\simeq\;
\FM_{|S|} \times \FM_{n - |S| + 1}\,,
\end{equation}
cluster times residual.

The \emph{ordered bar complex} is
\begin{equation}\label{eq:bar-ord}
\barB^{\mathrm{ord}}_X(\cA)
\;=\;
\bigoplus_{n \ge 1}
\bigl(s^{-1}\bar\cA\bigr)^{\otimes n}
\;\otimes\;
\Gamma\bigl(
\FM_n(X),\;
\Omega^{n-1}_{\log}
\bigr),
\end{equation}
with no $\Sigma_n$-quotient.  Here
$\bar\cA = \ker(\varepsilon \colon \cA \to \omega_X)$
is the augmentation ideal, $s^{-1}$ denotes desuspension
($|s^{-1}a| = |a| - 1$), and the cohomological convention
$|d| = +1$ is used throughout.

The bar differential decomposes:
\begin{equation}\label{eq:bar-diff}
d_{\barB}
\;=\;
d_{\mathrm{int}} + d_{\mathrm{res}} + d_{\mathrm{dR}}\,.
\end{equation}
The internal differential $d_{\mathrm{int}}$ acts on each
tensor factor; the de~Rham differential $d_{\mathrm{dR}}$ acts
on logarithmic forms on~$\FM_n(X)$; and the residue differential
$d_{\mathrm{res}}$ extracts OPE data at collision divisors.
For a codimension-one face $D_{\{i,j\}}$ corresponding to the
collision $z_i \to z_j$, the residue differential extracts
the $p$-th Fourier mode $a_{(p)}b$ of the OPE, where $p$ is
the pole order read from the form-theoretic coefficient of
$\eta_{ij}$.  The bar differential does not merely detect
singularities; it extracts them, mode by mode, through the
logarithmic kernel~\eqref{eq:eta}.


% ====================================================================
\subsection{Why $d^2 = 0$}
\label{ssec:d-squared}
% ====================================================================

Expand $(d_{\mathrm{int}} + d_{\mathrm{res}}
+ d_{\mathrm{dR}})^2$ into nine terms.  Three vanish
individually: $d_{\mathrm{int}}^2 = 0$ (the underlying cochain
complex), $d_{\mathrm{dR}}^2 = 0$ (de~Rham), and
$d_{\mathrm{int}}\,d_{\mathrm{dR}}
+ d_{\mathrm{dR}}\,d_{\mathrm{int}} = 0$ (Leibniz).  The
remaining terms cancel pairwise, governed by index overlap:

\begin{enumerate}[(a)]
\item \emph{Disjoint supports.}
When $\{i,j\} \cap \{k,l\} = \emptyset$, the faces
$D_{\{i,j\}}$ and $D_{\{k,l\}}$ are transverse
in~$\FM_n(X)$.  The residue operators commute:
$d_{\mathrm{res}}^{(ij)}\,d_{\mathrm{res}}^{(kl)}
+ d_{\mathrm{res}}^{(kl)}\,d_{\mathrm{res}}^{(ij)} = 0$.

\item \emph{Shared index.}
When $\{i,j\} \cap \{k,l\} = \{j = k\}$, the triple
collision $z_i = z_j = z_l$ sits in a codimension-two
stratum.  The form-theoretic contribution involves
$\eta_{ij} \wedge \eta_{jl}$.  The Arnold
relation~\eqref{eq:arnold} identifies this with a cyclic
sum over the three ordered approaches to the triple-collision
point.  On the algebraic side, the three terms produce
$(a_{i(m)}a_j)_{(n)}a_l$,
$(a_{j(m)}a_l)_{(n)}a_i$,
$(a_{l(m)}a_i)_{(n)}a_j$
with appropriate pole orders, whose sum vanishes by the
Borcherds identity~\eqref{eq:borcherds}.

The Arnold relation kills the form; the Borcherds identity
kills the algebra.  Both are needed, and both are consumed.

\item \emph{Same pair.}
$d_{\mathrm{res}}^{(ij)}\,d_{\mathrm{res}}^{(ij)} = 0$
by the sign rule on desuspensions: extracting a residue
twice at the same diagonal vanishes for degree reasons.
\end{enumerate}

The result: $d_{\barB}^2 = 0$ if and only if the
$n$-th products satisfy~\eqref{eq:borcherds}.  The bar
construction accepts exactly chiral algebras.


% ====================================================================
\subsection{Two coproducts}
\label{ssec:coproducts}
% ====================================================================

The ordered bar $\barB^{\mathrm{ord}}(\cA)$ carries the
\emph{deconcatenation coproduct}:
\begin{equation}\label{eq:deconc}
\Delta[s^{-1}a_1 | \cdots | s^{-1}a_n]
\;=\;
\sum_{i=0}^{n}
[s^{-1}a_1 | \cdots | s^{-1}a_i]
\;\otimes\;
[s^{-1}a_{i+1} | \cdots | s^{-1}a_n]\,,
\end{equation}
splitting an ordered sequence at every consecutive position:
$n + 1$ terms.  This is coassociative but not cocommutative.
It makes $\barB^{\mathrm{ord}}(\cA)$ a cofree tensor
coalgebra $T^c(s^{-1}\bar\cA)$~\cite{LV12}: the
$E_1$-coalgebra reflecting the linear ordering of points.

The $\Sigma_n$-coinvariant quotient
\begin{equation}\label{eq:bar-sym}
\barB^{\Sigma}_X(\cA)
\;=\;
\bigoplus_{n \ge 1}
\bigl(s^{-1}\bar\cA\bigr)^{\otimes n}_{\Sigma_n}
\;\otimes\;
\Gamma\bigl(
\FM_n(X),\;
\Omega^{n-1}_{\log}
\bigr)
\end{equation}
carries the \emph{coshuffle coproduct}, summing over all $2^n$
bipartitions of an unordered set: coassociative and
cocommutative.  This is the factorization coalgebra
of Beilinson--Drinfeld on~$\Ran(X)$~\cite{BD04}, the
$E_\infty$-coalgebra.

The ordered bar is the natural primitive.  The symmetric bar
is a quotient that loses data: the non-cocommutative structure
of the ordered bar carries the spectral $R$-matrix, the
Knizhnik--Zamolodchikov associator, and the full quantum group
deformation, none of which survive the passage to
$\Sigma_n$-coinvariants.

\begin{remark}[Six bar complexes]\label{rem:six-bars}
Six bar-type objects coexist.  The Francis--Gaitsgory
bar~\cite{FG12} $B^{\mathrm{Lie}}(\cA)$ is the cofree
coLie coalgebra $\mathrm{Lie}^c(s^{-1}\bar\cA)$, capturing
zeroth-pole data; it embeds into~$\barB^\Sigma$ via the
inclusion $\mathrm{Lie}^c \hookrightarrow \Sym^c$
(a quasi-isomorphism in characteristic zero, but a categorical
distinction).  The genus-graded bar $B^{(g)}(\cA)$ extends
$\barB^\Sigma$ to families of curves over
$\overline{\cM}_g$, where the differential acquires curvature
$d_{\mathrm{fib}}^2 = \kappa(\cA) \cdot \omega_g$
(see~\S\ref{ssec:genus}).
The modular bar $B^{\mathrm{mod}}(\cA)$ restores $D^2 = 0$
by completing over the Feynman transform of the commutative
modular operad.  The $E_1$-modular bar
$B^{E_1\text{-}\mathrm{mod}}(\cA)$, an algebra over the
Feynman transform of the associative modular operad, is the
primitive among these; the modular bar is its
$\Sigma_n$-coinvariant image.  No two of these six objects
are isomorphic in general.
\end{remark}


% ====================================================================
\subsection{Three operations on the bar}
\label{ssec:three-operations}
% ====================================================================

The bar construction is a categorical logarithm:
$B(\cA \otimes \cA') \simeq B(\cA) \oplus B(\cA')$.
The logarithmic form
$\eta_{ij} = d\log(z_i - z_j)$ is the integral kernel;
the Arnold relation is the cocycle condition for~$\log$.

Three distinct functors act on~$B(\cA)$ and produce three
distinct outputs:
\begin{equation}\label{eq:three-functors}
\begin{array}{lcll}
\Omega(B(\cA)) &\simeq& \cA
  & \text{(cobar inversion: recovers the original)},\\[3pt]
\mathbb{D}_{\Ran}(B(\cA)) &\simeq& \cA^!_\infty
  & \text{(Verdier duality: the homotopy Koszul dual)},
  \\[3pt]
(H^*(B(\cA)))^\vee &=& \cA^!
  & \text{(linear duality of bar cohomology: the strict dual)}.
\end{array}
\end{equation}
Cobar inversion recovers the original algebra~$\cA$ itself;
it is a round-trip, not a duality.  Verdier duality on
$\Ran(X)$ converts the bar coalgebra into a factorization
algebra~$\cA^!_\infty$, the homotopy Koszul dual retaining
full chain-level data.  That its cohomology recovers the
strict Koszul dual~$\cA^!$ on the Koszul locus is the content
of the Verdier intertwining theorem, not a tautology.
Confusing any two of~\eqref{eq:three-functors} is a
category error.


% ====================================================================
\subsection{The Heisenberg algebra at bar degree two}
\label{ssec:heisenberg}
% ====================================================================

The Heisenberg chiral algebra $\cH_k$ is generated by a
single field~$\alpha$ of conformal weight~$1$ with OPE
\begin{equation}\label{eq:heis-ope}
\alpha(z)\,\alpha(w)
\;\sim\;
\frac{k}{(z-w)^2}\,.
\end{equation}
Only the second-order pole is present:
$\alpha_{(1)}\alpha = k \cdot \mathbf{1}$ and
$\alpha_{(n)}\alpha = 0$ for $n \neq 1$.  No first-order
pole: $\alpha_{(0)}\alpha = 0$, no Lie bracket.
All triple products vanish since
$\mathbf{1}_{(n)} = 0$ for $n \ge 0$.

The bar complex at tensor degree~$2$ gives
\begin{equation}\label{eq:heis-bar2}
d_{\mathrm{res}}
(s^{-1}\alpha \otimes s^{-1}\alpha \otimes \eta_{12})
\;=\;
\alpha_{(1)}\alpha
\;=\;
k \cdot \mathbf{1}\,.
\end{equation}
The level~$k$ is visible already at genus~$0$.  At tensor
degree $n \ge 3$: every residue at a collision stratum
$D_{\{i,j\}}$ produces $\alpha_{(1)}\alpha = k \cdot \mathbf{1}$
on the pair~$(i,j)$; since $\mathbf{1}_{(p)}\alpha = 0$ for
$p \ge 0$, no further contractions occur.  All structure is
captured at tensor degree~$2$: the bar complex is quadratic.
The Heisenberg algebra is chirally Koszul.

The ordered bar carries the collision residue
\begin{equation}\label{eq:heis-r}
r(z) \;=\; \frac{k}{z}\,.
\end{equation}
The second-order OPE pole $k/(z-w)^2$ drops by one order under
the logarithmic kernel $d\log(z-w)$: the bar differential
absorbs one power of $(z-w)$, producing the first-order
$r$-matrix pole.  Applying the averaging map:
\begin{equation}\label{eq:heis-kappa}
\mathrm{av}\bigl(r(z)\bigr)
\;=\;
k
\;=\;
\kappa(\cH_k)\,.
\end{equation}
For the Heisenberg algebra, the $R$-matrix \emph{is}
the modular characteristic: $r(z) = k/z$ is scalar, so
averaging loses nothing.

The deficiency: no first-order pole, no Lie bracket, no
matrix structure in the collision residue, no quantum group.
Everything that makes representation theory nontrivial is
absent.  The simplest algebra where the $R$-matrix is
matrix-valued is the first algebra with a first-order pole:
the affine Kac--Moody algebra.

\begin{example}[Kac--Moody collision residue]
\label{ex:km}
The affine Kac--Moody algebra $\widehat{\fg}_k$ at
level~$k$ is generated by currents $J^a$,
$a = 1,\ldots,\dim\fg$, with OPE
\[
J^a(z)\,J^b(w)
\;\sim\;
\frac{f^{ab}{}_{c}\,J^c(w)}{z - w}
\;+\;
\frac{k\,\kappa^{ab}}{(z - w)^2}\,.
\]
Both poles are present.  The ordered bar carries the collision
residue
\begin{equation}\label{eq:km-r}
r(z)
\;=\;
\frac{k\,\Omega}{z}\,,
\qquad
\Omega = \sum_a J^a \otimes J_a
\;\in\;
\fg \otimes \fg\,,
\end{equation}
where $\Omega$ is the Casimir tensor and $k$ is the level.
At $k = 0$ the $r$-matrix vanishes: no level, no collision
data.  The averaging map collapses $k\,\Omega$ to a scalar:
\begin{equation}\label{eq:km-kappa}
\mathrm{av}(k\,\Omega/z)
\;=\;
\frac{\dim(\fg) \cdot (k + h^\vee)}{2h^\vee}
\;=\;
\kappa(\widehat{\fg}_k)\,,
\end{equation}
where $h^\vee$ is the dual Coxeter number.  The matrix
structure of~$\Omega$, erased by $\Sigma_2$-coinvariance,
encodes the representation theory of~$\fg$: it is the datum
that builds the Yangian~$Y(\fg)$.
\end{example}


% ====================================================================
\subsection{Genus~$0$ and genus $\ge 1$}
\label{ssec:genus}
% ====================================================================

Everything above takes place at genus~$0$: the propagator
$\eta_{12} = d\log(z_1 - z_2)$ is globally defined
on~$\mathbb{P}^1$, and the Arnold relation holds without
correction.  On a curve of genus $g \ge 1$, the function
$z_1 - z_2$ has no global meaning: the curve is compact,
and any meromorphic function with a single simple zero must
also have a simple pole.

The replacement is the \emph{prime form} $E(z_1, z_2)$, a
section of $K^{-1/2} \boxtimes K^{-1/2}$ vanishing to first
order along the diagonal with no other zeros.  The genus-$g$
propagator
\begin{equation}\label{eq:genus-g-prop}
\eta^{(g)}_{12}
\;=\;
d\log E(z_1, z_2)
+ \text{(period correction)}
\end{equation}
is single-valued but no longer holomorphic; the period
correction couples to the Hodge bundle
$\mathbb{E} \to \overline{\cM}_g$ through the period
matrix.  The propagator $d\log E(z_1,z_2)$ has conformal
weight~$1$ in both variables, regardless of the conformal
weight of the field being sewed: all edges of the genus-$g$
graph sum use the standard Hodge bundle
$\mathbb{E}_1 = R^0\pi_*\omega_\pi$.

The Arnold relation acquires a defect, and the fibrewise bar
differential satisfies
\begin{equation}\label{eq:curvature}
d_{\mathrm{fib}}^{\,2}
\;=\;
\kappa(\cA) \cdot \omega_g\,,
\end{equation}
where $\omega_g$ is the Arakelov form on the fibre of
$\pi \colon \mathcal{C}_g \to \overline{\cM}_g$ and
$\kappa(\cA) \in \Bbbk$ is the \emph{modular
characteristic}, determined entirely by the leading OPE
singularity.  Explicit values:
\begin{equation}\label{eq:kappa-values}
\kappa(\cH_k) = k\,,
\qquad
\kappa(\widehat{\fg}_k)
= \frac{\dim(\fg) \cdot (k + h^\vee)}{2h^\vee}\,,
\qquad
\kappa(\Vir_c)
= \frac{c}{2}\,.
\end{equation}
Period integrals restore nilpotence: the Gauss--Manin
connection absorbs the fibrewise curvature, and the total
differential satisfies $D_g^2 = 0$.

Genus~$0$ is uncurved algebra.  Genus $\ge 1$ is curved
geometry.  The entire genus tower is controlled by a single
scalar: the curvature of the bar complex on a family of
curves.

\begin{remark}[Recovering the classical bar]
\label{rem:classical}
When $X = \Spec\Bbbk$ is a point, configuration spaces
collapse, logarithmic forms vanish, and the chiral bar complex
recovers the classical bar construction of an associative
algebra~$A$.  This recovery requires specifying the retraction
data: vertex algebras live on the formal disk~$D$, and the
passage from~$D$ to a bare point discards the thickening
(completion, growth conditions) on which the
$\mathcal{D}$-module structure depends.
A quadratic algebra $A = T(V)/(R)$ has quadratic
dual $A^! = T(V^*)/(R^\perp)$; when~$A$ is Koszul, the bar
cohomology concentrates in diagonal degrees and
$H^*(B(A)) \cong (A^!)^\vee$.  The operadic instances
$\mathrm{Com}^! = \mathrm{Lie}$ and
$\Sym^! = \Lambda$ recover the Chevalley--Eilenberg complex
as a restricted bar construction.
\end{remark}


% ====================================================================
\section{Curvature and the genus tower}
\label{sec:curvature}
% ====================================================================

At genus~$0$ the bar complex is flat: $d^2 = 0$ is an
algebraic identity, equivalent to the Borcherds axiom.
On a family of curves of genus $g \ge 1$, this nilpotence
fails.  The failure is controlled by a single scalar.

% ====================================================================
\subsection{The prime form and the period defect}
\label{ssec:prime-form}
% ====================================================================

On a compact Riemann surface~$\Sigma_g$ of genus~$g \ge 1$,
the function $z_1 - z_2$ has no global meaning: a
meromorphic function with a single simple zero must have a
simple pole.  The \emph{prime form} $E(z_1, z_2)$, a
section of $K^{-1/2} \boxtimes K^{-1/2}$ vanishing to
first order along the diagonal, replaces it.  The
genus-$g$ propagator
\begin{equation}\label{eq:genus-g-propagator}
\eta^{(g)}_{12}
\;=\;
d\log E(z_1, z_2)
+ \pi \sum_{\alpha,\beta = 1}^{g}
(\operatorname{Im}\Omega)^{-1}_{\alpha\beta}
\,\omega_\alpha(z_1)\,\bar\omega_\beta(z_2)
\end{equation}
is single-valued; the period correction couples to the
Hodge bundle $\mathbb{E} = R^0\pi_*\omega_\pi$ through
the period matrix~$\Omega$.

Three properties of~\eqref{eq:genus-g-propagator} constrain
the genus tower.  First, the logarithmic derivative
$d\log E$ has conformal weight~$1$ in both variables,
regardless of the conformal weight of the field being
sewed: every edge of every graph sum at every genus uses
the standard Hodge bundle~$\mathbb{E}_1$.  Second, the
period correction is harmonic in each variable and involves
$(\operatorname{Im}\Omega)^{-1}$, the Arakelov Green
function on~$\Sigma_g$.  Third, the Arnold relation acquires
a defect proportional to the Arakelov form, and this defect
propagates into the bar differential.

% ====================================================================
\subsection{Curvature}
\label{ssec:curvature}
% ====================================================================

Let $\pi \colon \mathcal{C}_g \to \overline{\cM}_g$ be the
universal curve.  On each fibre~$\Sigma_g$, the propagator
$\eta^{(g)}$ is single-valued but no longer closed: the
period correction has a $\bar\partial$-component, and the
Arnold relation fails.  The fibrewise bar differential
satisfies
\begin{equation}\label{eq:d-squared-curvature}
d_{\mathrm{fib}}^{\,2}
\;=\;
\kappa(\cA) \cdot \omega_g\,,
\end{equation}
where $\omega_g$ is the Arakelov $(1,1)$-form on the fibre
and $\kappa(\cA) \in \Bbbk$ is the \emph{modular
characteristic}.  The modular characteristic is intrinsic to
the chiral algebra; it is determined by the leading OPE
singularity and is visible already at genus~$0$ as the
$\Sigma_2$-coinvariant of the collision residue:
$\kappa(\cA) = \mathrm{av}(r(z))$.

The failure of $d^2 = 0$ at genus~$g \ge 1$ is not a defect
of the construction; it is the construction detecting global
geometry.  Period integrals restore nilpotence: the
Gauss--Manin connection absorbs the fibrewise curvature, and
the total differential satisfies $D_g^2 = 0$.  The
curvature~\eqref{eq:d-squared-curvature} is the
infinitesimal avatar of this restoration.

% ====================================================================
\subsection{The genus expansion}
\label{ssec:genus-expansion}
% ====================================================================

For a modular Koszul chiral algebra~$\cA$ with a single
generator (or more generally with all generators of the
same conformal weight), the genus-$g$ obstruction class
$\mathrm{obs}_g(\cA)$ is proportional to the top Chern class
of the Hodge bundle:
\begin{equation}\label{eq:genus-universality}
\mathrm{obs}_g(\cA)
\;=\;
\kappa(\cA) \cdot \lambda_g\,,
\qquad g \ge 1\,.
\end{equation}
The free energy $F_g(\cA) = \kappa(\cA) \cdot
\lambda_g^{\mathrm{FP}}$, where
\begin{equation}\label{eq:faber-pandharipande}
\lambda_g^{\mathrm{FP}}
\;=\;
\frac{2^{2g-1} - 1}{2^{2g-1}}
\cdot
\frac{|B_{2g}|}{(2g)!}
\end{equation}
is the Faber--Pandharipande coefficient~\cite{FP00} ($B_{2g}$ the
$2g$-th Bernoulli number), packages the entire genus tower
into the Hirzebruch $\hat{A}$-class:
\begin{equation}\label{eq:ahat-generating}
\sum_{g=1}^{\infty}
F_g(\cA)\,\hbar^{2g-2}
\;=\;
\frac{\kappa(\cA)}{\hbar^2}
\bigl[\hat{A}(i\hbar) - 1\bigr]\,,
\qquad
\hat{A}(x) = \frac{x/2}{\sinh(x/2)}\,.
\end{equation}
At genus~$1$: $\lambda_1^{\mathrm{FP}} = 1/24$.  At
genus~$2$: $\lambda_2^{\mathrm{FP}} = 7/5760$.  The
universal ratio $F_2/F_1 = 7/240$ is independent
of~$\cA$.  The series converges for
$|\hbar| < 2\pi$, the first zero of $\sin(\hbar/2)$.

The genus-$1$ identity
$\mathrm{obs}_1 = \kappa \cdot \lambda_1$
is unconditional: it holds for all modular Koszul chiral
algebras, regardless of the number or weights of their
generators.  At genus $g \ge 2$, the story is more
delicate.

\begin{remark}[Multi-weight correction]
\label{rem:multi-weight}
For a chiral algebra with generators of distinct conformal
weights $h_1, \ldots, h_r$, the genus-$g$ free energy
decomposes as
\begin{equation}\label{eq:multi-weight-decomposition}
F_g(\cA)
\;=\;
\kappa(\cA) \cdot \lambda_g^{\mathrm{FP}}
\;+\;
\delta F_g^{\mathrm{cross}}(\cA)\,,
\end{equation}
where $\delta F_g^{\mathrm{cross}}$ is a graph sum over
mixed-channel boundary graphs of~$\overline{\cM}_{g,0}$.
At genus~$1$, $\delta F_1^{\mathrm{cross}} = 0$
unconditionally.  For uniform-weight algebras
($h_1 = \cdots = h_r$),
$\delta F_g^{\mathrm{cross}} = 0$ at all genera.  For
multi-weight algebras at $g \ge 2$, the correction is
generically nonzero: for $\cW_3$ at genus~$2$,
$\delta F_2^{\mathrm{cross}}(\cW_3)
= (c + 204)/(16c) > 0$ for all $c > 0$.
\end{remark}

% ====================================================================
\subsection{The modular characteristic}
\label{ssec:kappa}
% ====================================================================

The modular characteristic $\kappa(\cA)$ is the single
scalar controlling the genus tower.  It is computed from
genus-$0$ data alone: the leading coefficient of the OPE,
extracted through the bar construction.

\begin{table}[ht]
\centering
\caption{Modular characteristic for the standard families}
\label{tab:kappa}
\renewcommand{\arraystretch}{1.6}
\begin{tabular}{lcc}
\toprule
\textbf{Algebra $\cA$}
  & $\boldsymbol{\kappa(\cA)}$
  & \textbf{Shadow depth class} \\
\midrule
Heisenberg $\cH_k$
  & $k$
  & $\mathbf{G}$ \\
Affine Kac--Moody $\widehat{\fg}_k$
  & $\dfrac{\dim(\fg) \cdot (k + h^\vee)}{2h^\vee}$
  & $\mathbf{L}$ \\[4pt]
$\beta\gamma$ system
  & $c/2$
  & $\mathbf{C}$ \\
Virasoro $\Vir_c$
  & $c/2$
  & $\mathbf{M}$ \\
$\cW_3^k(\mathfrak{sl}_3)$
  & $5c/6$
  & $\mathbf{M}$ \\
$\cW_N^k(\mathfrak{sl}_N)$
  & $\displaystyle c \cdot \sum_{j=2}^{N}\frac{1}{j}$
  & $\mathbf{M}$ \\[4pt]
\bottomrule
\end{tabular}
\end{table}

The modular characteristic is additive under independent
sums ($\kappa(\cA \oplus \cA') = \kappa(\cA) + \kappa(\cA')$
when the mixed OPE vanishes) and constrained by Koszul
duality: $\kappa(\cA) + \kappa(\cA^!) = 0$ for
Kac--Moody and free-field algebras;
$\kappa(\Vir_c) + \kappa(\Vir_{26-c}) = 13$ for Virasoro.
In general, $\kappa(\cA) + \kappa(\cA^!) = \rho \cdot K$
for $\cW$-algebras, where $\rho$ is the
anomaly ratio and $K$ the complementarity constant.

At the self-dual point ($c = 13$ for Virasoro), the
complementarity sum is symmetric:
$\kappa(\Vir_{13}) = 13/2 = \kappa(\Vir_{13}^!)$.
At the critical level ($k = -h^\vee$ for Kac--Moody),
$\kappa = 0$ and the bar complex is flat at all genera:
the content migrates from curvature to bar cohomology,
where it appears as the Feigin--Frenkel centre~\cite{FF92}.


% ====================================================================
\section{The five main theorems}
\label{sec:five-theorems}
% ====================================================================

The bar construction of a chiral algebra on a curve produces
a factorization coalgebra on the Ran space of the
curve~\cite{BD04,CG17}.
The five main theorems describe what this coalgebra
determines: the Koszul dual, the original algebra,
the complementarity decomposition, the genus tower, and
the chiral Hochschild cohomology.  Each theorem is a
projection of the universal Maurer--Cartan element
$\Theta_\cA$ (Section~\ref{sec:mc}).

% ====================================================================
\subsection{Theorem A: Verdier intertwining}
\label{ssec:thm-a}
% ====================================================================

\begin{theorem}[A]\label{thm:A}
Let\/ $\cA$ be a chiral algebra on a smooth projective
curve~$X$.  There is an adjunction
\[
B
\;\colon\;
\mathrm{Alg}_{\mathrm{aug}}(\mathrm{Fact}(X))
\;\rightleftarrows\;
\mathrm{CoAlg}_{\mathrm{conil}}(\mathrm{Fact}(X))
\;\colon\;
\Omega
\]
between augmented factorization algebras and conilpotent
factorization coalgebras, with unit and counit the
universal twisting morphisms.  Verdier duality on\/
$\Ran(X)$ intertwines the bar of\/~$\cA$ with the bar
of its Koszul dual:
\begin{equation}\label{eq:verdier}
\mathbb{D}_{\Ran}\bigl(B(\cA)\bigr)
\;\simeq\;
B(\cA^!)\,,
\end{equation}
as factorization coalgebras, or equivalently
$\mathbb{D}_{\Ran}(B(\cA)) \simeq \cA^!_\infty$
as the homotopy Koszul dual algebra.
\end{theorem}

The adjunction is formal: it requires only the augmented
factorization-algebra structure and the duality on~$\Ran(X)$.
The intertwining~\eqref{eq:verdier} is a deeper statement:
it says that the Verdier dual of the bar \emph{coalgebra}
is the bar of the Koszul dual \emph{algebra}.  The
identification of the homotopy Koszul dual $\cA^!_\infty$
(retaining full chain-level data) with the strict Koszul
dual $\cA^! = (H^*(B(\cA)))^\vee$ (retaining only
cohomology) is the content of Theorem~A on the Koszul
locus, not a tautology.

% ====================================================================
\subsection{Theorem B: bar-cobar inversion}
\label{ssec:thm-b}
% ====================================================================

\begin{theorem}[B]\label{thm:B}
On the Koszul locus, the cobar of the bar recovers the
original algebra:
\begin{equation}\label{eq:inversion}
\Omega\bigl(B(\cA)\bigr)
\;\xrightarrow{\;\sim\;}
\;\cA\,.
\end{equation}
The counit of the bar-cobar adjunction is a
quasi-isomorphism of augmented factorization algebras.
\end{theorem}

Bar-cobar inversion is a round-trip, not a duality.
It does not produce the Koszul dual (that is Verdier
duality, Theorem~A); it does not produce the bulk
algebra (that is the chiral derived centre, Theorem~H).
It recovers~$\cA$ itself.

The proof proceeds by constructing a contracting homotopy
for the counit, using the PBW filtration as a control
parameter: the associated graded of the cobar complex
coincides with a cofree resolution of the PBW-graded
algebra, and the filtration spectral sequence degenerates
at~$E_2$.  Convergence of the spectral sequence requires
the Koszul hypothesis; for non-Koszul algebras, bar-cobar
inversion requires completion.

% ====================================================================
\subsection{Theorem C: complementarity}
\label{ssec:thm-c}
% ====================================================================

\begin{theorem}[C]\label{thm:C}
For a Koszul pair $(\cA, \cA^!)$ on~$X$ and $g \ge 2$, the
complementarity complex
\[
\mathbf{C}_g(\cA)
\;:=\;
R\Gamma\bigl(\overline{\cM}_g,\, \mathcal{Z}_\cA\bigr)
\]
decomposes into homotopy eigenspaces of the Verdier
involution~$\sigma$:
\begin{equation}\label{eq:complementarity}
\mathbf{C}_g
\;\simeq\;
\mathbf{Q}_g(\cA) \;\oplus\; \mathbf{Q}_g(\cA^!)\,.
\end{equation}
The summands $\mathbf{Q}_g(\cA) =
\operatorname{fib}(\sigma - \mathrm{id})$
and\/ $\mathbf{Q}_g(\cA^!) =
\operatorname{fib}(\sigma + \mathrm{id})$ are
Lagrangian for the $(-1)$-shifted symplectic pairing
on\/~$\mathbf{C}_g$.
\end{theorem}

The eigenspace decomposition (C1) is unconditional: it
holds for all modular Koszul algebras at all genera.  The
scalar identity $F_g = \kappa \cdot \lambda_g^{\mathrm{FP}}$
on each summand (C2) requires the uniform-weight hypothesis
at $g \ge 2$.  The Lagrangian property, a
$(-1)$-shifted symplectic statement in the sense of
Pantev--To\"en--Vaqui\'e--Vezzosi~\cite{PTVV13}, is conditional on
perfectness and nondegeneracy of the cyclic pairing.

% ====================================================================
\subsection{Theorem D: the modular characteristic}
\label{ssec:thm-d}
% ====================================================================

\begin{theorem}[D]\label{thm:D}
For a uniform-weight modular Koszul chiral algebra\/~$\cA$,
the genus-$g$ obstruction class satisfies
$\mathrm{obs}_g(\cA) = \kappa(\cA) \cdot \lambda_g$
for all\/ $g \ge 1$, where $\kappa(\cA)$ is a universal,
additive, duality-constrained scalar determined by the
genus-$0$ data of\/~$\cA$.  The generating function is
the Hirzebruch $\hat{A}$-class
\textup{(}\eqref{eq:ahat-generating}\textup{)}.
At genus\/~$1$, the identity is unconditional: it holds
for all modular Koszul algebras regardless of weight
structure.
\end{theorem}

The universality is the point: no matter how complicated
the OPE structure, the genus tower is controlled by a
single number.  Theorem~D is the all-genera identity on the
uniform-weight lane; at genus~$1$ it is unconditional
(Remark~\ref{rem:multi-weight}).

% ====================================================================
\subsection{Theorem H: chiral Hochschild cohomology}
\label{ssec:thm-h}
% ====================================================================

\begin{theorem}[H]\label{thm:H}
For a chirally Koszul algebra~$\cA$, the chiral Hochschild
cohomology
$\mathrm{ChirHoch}^*(\cA)
= H^*(C^{\bullet}_{\mathrm{ch}}(\cA, \cA), \delta)$
is polynomial with generators concentrated in degrees
$\{0, 1, 2\}$:
\[
\dim_\Bbbk \mathrm{ChirHoch}^*(\cA)
\;\le\; 4\,.
\]
\end{theorem}

The chiral derived centre
$Z^{\mathrm{der}}_{\mathrm{ch}}(\cA) =
C^{\bullet}_{\mathrm{ch}}(\cA, \cA)$ is the universal
bulk algebra: the algebra of operators in the interior
that act on the boundary algebra~$\cA$.  Theorem~H
constrains its cohomology.  For the Virasoro algebra,
the total dimension of $\mathrm{ChirHoch}^*(\Vir_c)$ is
at most~$4$; this is a finite-dimensionality statement,
not a polynomial-ring statement.

% ====================================================================
\subsection{The Koszulness meta-theorem}
\label{ssec:meta-theorem}
% ====================================================================

The hypothesis shared by Theorems~A--H is chiral
Koszulness: bar cohomology $H^*(B(\cA))$ concentrated in
bar degree~$1$.  The meta-theorem characterizes this
condition in twelve equivalent ways, ten of which are
unconditional, one conditional, one partial.

\begin{theorem}[Koszulness characterization]
\label{thm:koszul-meta}
For a chiral algebra~$\cA$ on~$X$ with PBW filtration,
conditions \textup{(i)--(x)} are equivalent.  Under
additional perfectness hypotheses, \textup{(xi)} is
equivalent.  Condition \textup{(xii)} implies
\textup{(x)}; the converse is open.
\begin{enumerate}[label=\textup{(\roman*)}]
\item $\cA$ is chirally Koszul.
\item The PBW spectral sequence collapses at\/~$E_2$.
\item The minimal $A_\infty$-model of $B(\cA)$
  is formal: $m_n = 0$ for $n \ge 3$.
\item Ext diagonal vanishing:
  $\operatorname{Ext}^{p,q}_\cA(\omega_X, \omega_X) = 0$
  for $p \ne q$.
\item The bar-cobar counit
  $\Omega(B(\cA)) \to \cA$ is a quasi-isomorphism.
\item The Barr--Beck--Lurie comparison is an equivalence.
\item Factorization homology
  $\int_{\Sigma_g}\!\cA$ is concentrated in degree~$0$
  for all~$g$.
\item $\mathrm{ChirHoch}^*(\cA)$ is polynomial in
  degrees $\{0, 1, 2\}$.
\item The Kac--Shapovalov determinant is nonzero in the
  bar-relevant range.
\item Restriction to every FM boundary stratum is
  acyclic outside degree~$0$.
\item[{\upshape (xi)}] Lagrangian criterion \textup{(conditional)}.
\item[{\upshape (xii)}] $\cD$-module purity
  \textup{(one-directional: (x) $\Rightarrow$ (xii))}.
\end{enumerate}
\end{theorem}

The circuit of unconditional equivalences runs through five
independent mathematical domains: PBW degeneration
(commutative algebra), $A_\infty$-formality (homotopical
algebra), Barr--Beck--Lurie (higher categories),
Kac--Shapovalov (representation theory), and FM boundary
acyclicity (algebraic topology).  That these five conditions
define the same class of algebras is the depth of the
meta-theorem.


% ====================================================================
\section{The universal Maurer--Cartan element
  and the shadow tower}
\label{sec:mc}
% ====================================================================

The five theorems of Section~\ref{sec:five-theorems} are
projections of a single algebraic object: the universal
Maurer--Cartan element $\Theta_\cA$ of the modular
convolution algebra~\cite{RNW19,Val16}.

% ====================================================================
\subsection{The bar-intrinsic construction}
\label{ssec:bar-intrinsic}
% ====================================================================

The genus-completed bar differential
$D_\cA = \sum_{g \ge 0} \hbar^g\,d^{(g)}_\cA$
satisfies $D_\cA^2 = 0$ by the codimension-$2$ face
cancellation on $\overline{\cM}_{g,n}$ (the same mechanism
as~\S\ref{ssec:d-squared}, extended to all genera by Mok's
log Fulton--MacPherson compactification~\cite{Mok25}).  Write
$d_0 = d^{(0)}_\cA$ for the genus-$0$ bar differential.

\begin{construction}[Bar-intrinsic MC element]
\label{constr:bar-intrinsic}
Define the positive-genus correction
\begin{equation}\label{eq:theta-def}
\Theta_\cA
\;:=\;
D_\cA - d_0
\;=\;
\sum_{g \ge 1}
\hbar^g\,d^{(g)}_\cA
\;\in\;
\prod_{g \ge 1}
\mathrm{CoDer}^{\mathrm{cyc}}
\bigl(B^{(g)}_X(\cA)\bigr)[1]\,.
\end{equation}
\end{construction}

\begin{theorem}\label{thm:theta-mc}
$\Theta_\cA$ is a Maurer--Cartan element of the modular
cyclic deformation complex:
\begin{equation}\label{eq:mc-equation}
d_0(\Theta_\cA)
+ \tfrac{1}{2}[\Theta_\cA, \Theta_\cA]
\;=\; 0\,.
\end{equation}
\end{theorem}

\begin{proof}[Proof sketch]
Expand $D_\cA^2 = (d_0 + \Theta_\cA)^2 = 0$.
Since $d_0^2 = 0$ (genus-$0$ bar flatness), the
remaining terms give
$d_0(\Theta_\cA) + \frac{1}{2}[\Theta_\cA, \Theta_\cA]
= 0$, which is the MC equation.  The bracket is the
convolution bracket on the modular cyclic deformation
complex, inherited from stable-curve gluing; its square-zero
property is a formal consequence of $d_0^2 = 0$.
\end{proof}

The construction requires no restriction to simple Lie
symmetry, no one-channel hypothesis, and no tautological-line
support condition.  It works for all modular Koszul chiral
algebras~\cite{MS24}.  The modular characteristic~$\kappa$, the cubic
shadow~$\mathfrak{C}$, and the quartic resonance
class~$\mathfrak{Q}$ are arity projections of this single
element:
\begin{equation}\label{eq:shadow-projections}
\kappa(\cA) = \pi_2(\Theta_\cA)\,,
\qquad
\mathfrak{C}(\cA) = \pi_3(\Theta_\cA)\,,
\qquad
\mathfrak{Q}(\cA) = \pi_4(\Theta_\cA)\,.
\end{equation}
Each arity projection carries more information than the
last; the full $\Theta_\cA$ carries all of it.

% ====================================================================
\subsection{The shadow obstruction tower}
\label{ssec:shadow-tower}
% ====================================================================

The arity-$r$ projection
$\mathrm{Sh}_r(\cA) := \pi_r(\Theta_\cA)$ is the $r$-th
shadow coefficient.  The sequence $S_2 = \kappa$,
$S_3 = \alpha$, $S_4$, $S_5, \ldots$ constitutes the
\emph{shadow obstruction tower}.  On any one-dimensional
primary slice~$L$ of the cyclic deformation complex, the
MC equation~\eqref{eq:mc-equation} imposes a quadratic
constraint.

\begin{theorem}[Riccati algebraicity]
\label{thm:riccati}
The weighted shadow generating function
$H(t) := \sum_{r \ge 2} r\,S_r\,t^r$
satisfies
\begin{equation}\label{eq:algebraic-relation}
H(t)^2 \;=\; t^4\,Q_L(t)\,,
\end{equation}
where
\begin{equation}\label{eq:shadow-metric}
Q_L(t)
\;=\;
(2\kappa + 3\alpha t)^2
\;+\;
2\Delta\,t^2\,,
\qquad
\Delta := 8\kappa S_4\,,
\end{equation}
is the \emph{shadow metric}: a quadratic polynomial in~$t$
determined by three genus-$0$ invariants
$(\kappa, \alpha, S_4)$.
The entire shadow obstruction tower on~$L$ is determined
by~$Q_L$.
\end{theorem}

The generating function $H(t) = t^2\sqrt{Q_L(t)}$ is
algebraic of degree~$2$.  The shadow obstruction tower
is neither a Taylor expansion that might diverge nor a
formal power series requiring analytic continuation: it is
the expansion of an explicit algebraic function.  The
MC equation, which a priori is an infinite system of coupled
nonlinear equations, reduces on each primary slice to a
single quadratic identity.

% ====================================================================
\subsection{The G/L/C/M classification}
\label{ssec:glcm}
% ====================================================================

The discriminant $\Delta = 8\kappa S_4$ of the shadow
metric forces a universal dichotomy.

\begin{definition}[Shadow depth classification]
\label{def:shadow-depth}
The \emph{shadow depth} of~$\cA$ is
$r_{\max}(\cA) := \sup\{r \ge 2 : S_r \ne 0\}$.
The shadow depth class is:
\begin{center}
\renewcommand{\arraystretch}{1.2}
\begin{tabular}{clcl}
\toprule
\textbf{Class} & \textbf{Name} &
  $\boldsymbol{r_{\max}}$ & \textbf{Mechanism} \\
\midrule
$\mathbf{G}$ & Gaussian &
  $2$ & $Q_L$ perfect square, $\alpha = 0$ \\
$\mathbf{L}$ & Lie &
  $3$ & $Q_L$ perfect square, $\alpha \ne 0$ \\
$\mathbf{C}$ & Contact &
  $4$ & stratum separation \\
$\mathbf{M}$ & Mixed &
  $\infty$ & $\Delta \ne 0$: $Q_L$ irreducible \\
\bottomrule
\end{tabular}
\end{center}
\end{definition}

The mechanism: when $\Delta = 0$, the shadow metric is a
perfect square $Q_L(t) = (2\kappa + 3\alpha t)^2$ and
$\sqrt{Q_L}$ is polynomial: the tower terminates.
If $\alpha = 0$ it terminates at arity~$2$ (class~$\mathbf{G}$:
Heisenberg); if $\alpha \ne 0$ at arity~$3$
(class~$\mathbf{L}$: affine Kac--Moody).  When
$\Delta \ne 0$, $\sqrt{Q_L}$ is irrational: the tower
runs forever (class~$\mathbf{M}$: Virasoro,
$\cW_N$).  Class~$\mathbf{C}$ ($\beta\gamma$)
escapes the dichotomy by stratum separation: the quartic
contact invariant lives on a charged stratum whose
self-bracket exits the complex by rank-one rigidity.

The four classes are a homotopy-invariant classification:
they coincide with the $L_\infty$-formality depth of the
transferred minimal model.  Class~$\mathbf{G}$ is
$L_\infty$-formal (all higher brackets vanish).
Class~$\mathbf{L}$ has exactly one nontrivial Massey
product.  Class~$\mathbf{C}$ has formality obstructions
through arity~$4$.  Class~$\mathbf{M}$ is intrinsically
non-formal: no finite truncation of the $L_\infty$
structure is quasi-isomorphic to a dg~Lie algebra.

Archetypes:
\begin{center}
\renewcommand{\arraystretch}{1.3}
\begin{tabular}{clll}
\toprule
\textbf{Class} & \textbf{Archetype}
  & $\boldsymbol{\kappa}$ & $\boldsymbol{\Delta}$ \\
\midrule
$\mathbf{G}$ & Heisenberg $\cH_k$
  & $k$ & $0$ \\
$\mathbf{L}$ & $\widehat{\fg}_k$
  & $\dim(\fg)(k{+}h^\vee)/(2h^\vee)$ & $0$ \\
$\mathbf{C}$ & $\beta\gamma$
  & $c/2$ & $0$ (stratum separation) \\
$\mathbf{M}$ & $\Vir_c$
  & $c/2$ & $40/(5c{+}22)$ \\
\bottomrule
\end{tabular}
\end{center}

For the Virasoro algebra, the quartic contact invariant is
$Q^{\mathrm{contact}}_{\Vir} = 10/[c(5c+22)]$ and the
critical discriminant is $\Delta_{\Vir} = 8 \cdot (c/2)
\cdot 10/[c(5c+22)] = 40/(5c+22)$.  Every shadow
coefficient $S_r$ for $r \ge 5$ is a rational function
of~$c$ with poles only at $c = 0$ and $c = -22/5$:
\begin{align}
S_5 &= -\frac{48}{c^2(5c{+}22)}\,,
  &
S_6 &= \frac{80(45c{+}193)}{3\,c^3(5c{+}22)^2}\,,
  &
S_7 &= -\frac{2880(15c{+}61)}{7\,c^4(5c{+}22)^2}\,.
\label{eq:shadow-coefficients}
\end{align}
These are the fingerprints of a class-$\mathbf{M}$ algebra:
the shadow tower is infinite, alternating in sign, and
controlled by an irrational algebraic function.

% ====================================================================
\subsection{The gauge-gravity dichotomy}
\label{ssec:gauge-gravity}
% ====================================================================

The classification has a physical interpretation.  At
$\kappa = 0$ the bar complex is flat at all genera:
the genus tower carries no curvature, the free energy
vanishes, and the content lives in bar cohomology.  This
is the regime of gauge theory: the Feigin--Frenkel centre
(the Langlands dual of the affine algebra at critical level)
is the genus-$0$ bar cohomology $H^0(B(\widehat{\fg}_{-h^\vee}))$.

At $\kappa \ne 0$ the bar complex is curved: every genus
contributes a nonzero obstruction class, the genus tower
is nontrivial, and the free energy is the $\hat{A}$-genus.
This is the regime of gravity: the bar curvature is the
perturbative genus expansion.

The Koszul duality $\cA \leftrightarrow \cA^!$ exchanges
gauge and gravity sectors.  For affine Kac--Moody at
generic level, $\kappa(\widehat{\fg}_k) \ne 0$ and the
algebra is curved; at the Koszul-dual critical level
$k' = -k - 2h^\vee$, $\kappa(\widehat{\fg}_{k'}) \ne 0$
as well, and both sides are gravitational.  The gauge regime
(critical level, $\kappa = 0$) is the fixed locus of a
one-parameter family.

For the Virasoro algebra, the Koszul dual is
$\Vir_c^! = \Vir_{26-c}$.  The gauge locus is $c = 0$
(topological: $\kappa = 0$); the gravitational locus is
$c \ne 0$.  The self-dual point $c = 13$
($\kappa = 13/2 = \kappa^!$) is the critical string.
The Clifford algebra of the complementarity pairing
degenerates at $c = 26$ ($\kappa(\Vir_0^!) = 0$):
the $26$-dimensional bosonic string is the critical point
where the Koszul dual becomes flat.

The dichotomy gauge/$\kappa{=}0$ vs
gravity/$\kappa{\ne}0$ is not a physical input: it is a
theorem about the bar complex.  The physical interpretation
follows from the identification of bar curvature with
perturbative genus expansions in quantum field theory.


% ====================================================================
\section{The collision residue and the seven faces}
\label{sec:seven-faces}
% ====================================================================

The bar differential extracts OPE data at collision divisors.
When two points collide, the residue of the bar-intrinsic MC
element $\Theta_\cA$ at the codimension-one stratum
$D_{\{1,2\}} \subset \FM_2(X)$ defines the \emph{collision
residue}:
\begin{equation}\label{eq:collision-residue}
r_\cA(z)
\;:=\;
\Res^{\mathrm{coll}}_{0,2}(\Theta_\cA)
\;\in\;
\End_\cA(2)[\![z^{-1}]\!]\,.
\end{equation}
The collision residue is the arity-$2$, genus-$0$ projection
of $\Theta_\cA$.  Its pole orders are one less than those of
the OPE: the logarithmic kernel $d\log(z-w)$ absorbs one
power of $(z - w)$.  For a chiral algebra with OPE poles of
maximal order~$p$, the collision residue has poles of
order~$p - 1$.

% ====================================================================
\subsection{Pole-order dichotomy}
\label{ssec:pole-dichotomy}
% ====================================================================

The maximal pole order of the collision residue separates the
standard landscape into two regimes:

\begin{center}
\renewcommand{\arraystretch}{1.3}
\begin{tabular}{clcl}
\toprule
\textbf{Pole order} & \textbf{Algebra}
  & \textbf{Collision} & \textbf{Class} \\
\midrule
$1$ & Heisenberg $\cH_k$ & $k/z$ & $\mathbf{G}$ \\
$1$ & Kac--Moody $\widehat{\fg}_k$ & $k\,\Omega/z$ & $\mathbf{L}$ \\
$3$ & Virasoro $\Vir_c$ & $(c/2)/z^3 + 2T/z$ & $\mathbf{M}$ \\
$2N{-}1$ & $\cW_N$ & poles through $z^{-(2N-1)}$
  & $\mathbf{M}$ \\
\bottomrule
\end{tabular}
\end{center}

\noindent
This table encodes a structural dichotomy: algebras whose
collision residue has at most a simple pole (classes~$\mathbf{G}$
and~$\mathbf{L}$) have finite shadow depth and
Swiss-cheese-formal bar complexes; algebras with higher-order
poles (class~$\mathbf{M}$) have infinite shadow towers and
genuinely non-formal $A_\infty$-structure.

The dichotomy is not a property of individual OPE coefficients
but of the collision residue as a whole: the $d\log$ absorption
converts the OPE's double pole (carrying the metric) and simple
pole (carrying the Lie bracket) into a single simple pole
(carrying the Casimir tensor) for Kac--Moody, while the
Virasoro OPE's quartic pole becomes a cubic collision pole,
producing an $r$-matrix with genuine spectral dependence.

% ====================================================================
\subsection{Seven faces of $r(z)$}
\label{ssec:seven-faces}
% ====================================================================

The collision residue is a single object viewed through seven
algebraic lenses.

\begin{theorem}[Seven-face identification]
\label{thm:seven-faces}
For a modular Koszul chiral algebra~$\cA$, the following seven
expressions are equal as elements of\/
$\End_\cA(2)[\![z^{-1}]\!]$:
\begin{enumerate}[label=\textup{(F\arabic*)}]
\item The bar collision residue:
  $r_\cA(z)
  = \Res^{\mathrm{coll}}_{0,2}(\Theta_\cA)$.
\item The spectral $R$-matrix of the ordered bar:
  $R(z) = \mathrm{id} + r(z)/z + \cdots$ solves the
  classical Yang--Baxter equation.
\item The $\lambda$-bracket of the Poisson vertex algebra:
  $r(z) = \int_0^\infty e^{-\lambda z}
  \{a_\lambda b\}\,d\lambda$
  \textup{(}Laplace transform\textup{)}.
\item The Knizhnik--Zamolodchikov connection on
  $\Conf_n(\Bbbk)$:
  $\omega_{\mathrm{KZ}}
  = \sum_{i < j} r_{ij}(z_i {-} z_j)\,d(z_i {-} z_j)$,
  whose flatness is equivalent to the CYBE for~$r$.
\item\label{it:drinfeld-face} The Drinfeld $r$-matrix
  \textup{(}for $\cA = \widehat{\fg}_k$\textup{)}:
  $r(z) = k\,\Omega/z$.
\item The Sklyanin bracket:
  $\{L_1(u), L_2(v)\}
  = [r(u - v), L_1(u) \otimes L_2(v)]$.
\item The Gaudin Hamiltonians:
  $H_i
  = \sum_{j \ne i} r(z_i - z_j)$
  are mutually commuting.
\end{enumerate}
Faces~\textup{(F1)--(F4)} hold for every modular Koszul
algebra.  Faces~\textup{(F5)--(F7)} specialize to affine
Kac--Moody; the general case is obtained by replacing
$k\,\Omega/z$ with $r_\cA(z)$.
\end{theorem}

The seven faces are not seven theorems but seven readings of
one equation.  The collision residue
$r(z) = \pi_2(\Theta_\cA)$ is the arity-$2$ projection of
the universal MC element; the seven faces are the seven
classical frameworks in which this projection has been
studied.  Drinfeld's $r$-matrix~\cite{Dri85}, the KZ
connection~\cite{FBZ04}, the Sklyanin
bracket~\cite{FBZ04}, and the Gaudin model~\cite{FF92}
all extract the same binary invariant.

The relationship to the modular characteristic is immediate:
\begin{equation}\label{eq:kappa-from-r}
\kappa(\cA)
\;=\;
\mathrm{av}\bigl(r(z)\bigr)\,.
\end{equation}
The averaging map $\mathrm{av}$ is the
$\Sigma_2$-coinvariant projection.  For Kac--Moody,
$\mathrm{av}(k\,\Omega/z) = \dim(\fg)(k + h^\vee)/(2h^\vee)$.
For Virasoro, $\mathrm{av}((c/2)/z^3 + 2T/z) = c/2$.

% ====================================================================
\subsection{DS reduction as a functor on collision residues}
\label{ssec:ds-reduction}
% ====================================================================

Drinfeld--Sokolov reduction converts affine Kac--Moody
algebras into $\cW$-algebras:
$\widehat{\fg}_k \xmapsto{\mathrm{DS}}
\cW_N^k(\fg)$.  At the level of the collision
residue, this passage raises the pole order:

\begin{center}
\renewcommand{\arraystretch}{1.2}
\begin{tabular}{lccl}
\toprule
& \textbf{OPE max.\ pole} & \textbf{Collision max.\ pole}
  & \textbf{Class} \\
\midrule
$\widehat{\mathfrak{sl}}_N$ & $2$ & $1$ & $\mathbf{L}$ \\
$\cW_N$ & $2N$ & $2N{-}1$ & $\mathbf{M}$ \\
\bottomrule
\end{tabular}
\end{center}

\noindent
The Sugawara construction generates a Virasoro field
$T(z)$ from the affine currents $J^a(z)$.  The OPE
$T(z)\,T(w)$ has a quartic pole; DS reduction extends
this mechanism to all generators, producing a
$\cW$-algebra whose maximal OPE pole grows with
$N$.  The shadow depth jumps from finite (class~$\mathbf{L}$)
to infinite (class~$\mathbf{M}$): DS reduction is the functor
that transports gauge-type algebras into gravitational-type
algebras.

The intertwining theorem controls the Koszul-dual behavior:
the bar-cobar adjunction commutes with principal DS
reduction.  For non-principal reductions~\cite{KRW03}
(hook-type in type $A$, proved; arbitrary nilpotent,
programme), the
intertwining holds on the hook-type corridor and is
conjectural in full generality.


% ====================================================================
\section{The Swiss-cheese realization}
\label{sec:swiss-cheese}
% ====================================================================

The bar complex of a chiral algebra on a curve $X$ has two
structures: a differential (from OPE residues on Fulton--MacPherson
compactifications of~$X$) and a coproduct (from interval-cutting
in a topological direction).  Together, these produce an algebra
over the Swiss-cheese operad $\mathrm{SC}^{\mathrm{ch,top}}$,
the two-coloured operad governing the interaction of holomorphic
and topological factorization.

% ====================================================================
\subsection{The two colours}
\label{ssec:two-colours}
% ====================================================================

The bar complex lives on $\FM_k(\Bbbk) \times \Conf_k(\mathbb{R})$,
the product of holomorphic and topological configuration spaces.
This product is the operadic fingerprint of the Swiss-cheese
operad:

\begin{center}
\renewcommand{\arraystretch}{1.3}
\begin{tabular}{lll}
\toprule
& \textbf{Closed colour} & \textbf{Open colour} \\
\midrule
\textbf{Space} & $\FM_k(\Bbbk)$ & $\Conf_k(\mathbb{R})$ \\
\textbf{Structure} & Bar differential $d_{\barB}$ &
  Deconcatenation $\Delta$ \\
\textbf{Physics} & Holomorphic factorization & Topological
  factorization \\
\textbf{Operadic type} & $E_\infty$ & $E_1$ \\
\textbf{Coalgebra} & $\Sym^c(s^{-1}\bar\cA)$ &
  $T^c(s^{-1}\bar\cA)$ \\
\textbf{Coproduct} & Coshuffle ($2^n$ terms) &
  Deconcatenation ($n+1$ terms) \\
\bottomrule
\end{tabular}
\end{center}

\noindent
The closed colour is the holomorphic factorization of
Section~\ref{sec:bar}: the bar differential extracts OPE
residues, produces $d^2 = 0$ at genus~$0$, and acquires
curvature $\kappa(\cA) \cdot \omega_g$ at higher genus.  The
open colour is the topological factorization: the
deconcatenation coproduct splits an ordered sequence at every
consecutive position, producing the cofree tensor coalgebra.

The directionality of the Swiss-cheese operad is strict:
\emph{no open inputs produce closed outputs}.  Bulk operators
restrict to boundary operators, but not conversely.  The open
colour is the $E_1$-direction; the passage to
$\Sigma_n$-coinvariants (the closed/symmetric bar) is a lossy
projection.  The five main theorems A--H are the invariants
that survive this projection.

% ====================================================================
\subsection{Three bar complexes}
\label{ssec:three-bars}
% ====================================================================

Three bar-type objects coexist, corresponding to three
cooperadic structures on the same underlying graded vector
space:

\begin{enumerate}[(a)]
\item $\barB^{\mathrm{ord}}(\cA) = T^c(s^{-1}\bar\cA)$:
  the \emph{ordered bar}.  Deconcatenation coproduct ($n + 1$
  terms), coassociative but not cocommutative.  Carries the
  $R$-matrix, the KZ associator, and the full Yangian
  deformation.  This is the $E_1$-coalgebra.

\item $\barB^{\Sigma}(\cA) = \Sym^c(s^{-1}\bar\cA)$:
  the \emph{symmetric bar}.  Coshuffle coproduct ($2^n$
  terms), cocommutative and coassociative.  This is the
  factorization coalgebra on $\Ran(X)$: the $E_\infty$-coalgebra
  of Beilinson--Drinfeld.

\item $B^{\mathrm{FG}}(\cA)
  = \mathrm{Lie}^c(s^{-1}\bar\cA)$: the \emph{Francis--Gaitsgory
  bar}.  CoLie cobracket, capturing only the zeroth-pole OPE
  data ($a_{(0)}b$, the chiral Lie bracket).  A
  quasi-isomorphism $B^{\mathrm{FG}} \hookrightarrow
  \barB^\Sigma$ holds in characteristic zero, but
  $B^{\mathrm{FG}}$ is invisible to the curvature $\kappa(\cA)$:
  the modular characteristic comes from the \emph{first} OPE
  pole, which the Francis--Gaitsgory bar does not see.
\end{enumerate}

The natural descent is
$\barB^{\mathrm{ord}} \xrightarrow{\mathrm{av}}
\barB^{\Sigma} \hookleftarrow B^{\mathrm{FG}}$:
the ordered bar maps to the symmetric bar by
$\Sigma_n$-coinvariance (a lossy projection), and the
Francis--Gaitsgory bar embeds into the symmetric bar
(a quasi-isomorphism in characteristic zero, losing all
higher-pole data).  The ordered bar is the natural
primitive; the other two are its images.

% ====================================================================
\subsection{PVA descent and the recognition theorem}
\label{ssec:pva}
% ====================================================================

The cohomology of a Swiss-cheese algebra carries a Poisson
vertex algebra structure: the commutative product descends
from the regular OPE and the Poisson bracket from the singular
OPE.  This is the \emph{PVA descent}: the classical shadow of
the full quantum Swiss-cheese structure.

The six PVA axioms (Leibniz, skew-symmetry, Jacobi, sesqui\-linearity,
$\partial$-bracket, and Wick) are proved geometrically from
configuration-space geometry: the exchange cylinder on
$\FM_2(\Bbbk)$ for skew-symmetry, the three-face Stokes
theorem on $\FM_3(\Bbbk)$ for Jacobi, and boundary
factorization on higher $\FM_k(\Bbbk)$ for the remaining
axioms.

The recognition theorem establishes the converse:

\begin{theorem}[Recognition]
\label{thm:recognition}
A factorization algebra on\/ $\Ran(X)$ equipped with Weiss
cosheaf descent data and compatible with the factorization
isomorphisms is a $C_*(\mathcal{W}(\mathrm{SC}^{\mathrm{ch,top}}))$-algebra.
\end{theorem}

The theorem converts the abstract factorization structure
into a concrete operadic algebra over the Boardman--Vogt
resolution of the Swiss-cheese operad.  It is proved by Weiss
cosheaf descent: the factorization pattern on $\Ran(X)$
determines the operadic composition uniquely.

The homotopy-Koszulity of $\mathrm{SC}^{\mathrm{ch,top}}$
is proved via Kontsevich formality and transfer from the
classical Swiss-cheese operad (Livernet).  This removes all
previously conditional hypotheses: bar-cobar Quillen
equivalence, filtered Koszul duality, and the identification
of line operators with dual modules all become unconditional.


% ====================================================================
\section{Three-dimensional HT QFT and gauge-gravity}
\label{sec:ht-qft}
% ====================================================================

A four-dimensional holomorphic-topological field theory on
$\Bbbk_z \times \mathbb{R}_t \times \mathbb{R}_{>0}$
restricts to a chiral algebra on the holomorphic boundary.
The bar complex of that boundary algebra classifies
twisting morphisms (couplings to the Koszul dual); the
chiral derived centre computes the universal bulk; bar-cobar
inversion recovers the boundary algebra itself.  The modular
convolution algebra organizes the full boundary-to-bulk
correspondence at all genera.

% ====================================================================
\subsection{The holographic datum}
\label{ssec:holographic-datum}
% ====================================================================

The complete data of a holographic system is a sextuple:

\begin{definition}[Holographic modular Koszul datum]
\label{def:holographic}
The \emph{holographic modular Koszul datum} of a chiral
algebra~$\cA$ is the sextuple
\begin{equation}\label{eq:holographic-datum}
\cH(\cA)
\;=\;
\bigl(\cA,\;
\cA^!,\;
\mathcal{C},\;
r(z),\;
\Theta_\cA,\;
\nabla^{\mathrm{hol}}\bigr)\,,
\end{equation}
where $\cA^! = (H^*(B(\cA)))^\vee$ is the Koszul dual,
$\mathcal{C} \simeq \cA^!\text{-}\mathsf{mod}$ is the
line-operator category, $r(z)$ is the collision residue,
$\Theta_\cA$ is the universal MC element, and
$\nabla^{\mathrm{hol}}_{g,n} = d -
\mathrm{Sh}_{g,n}(\Theta_\cA)$ is the modular shadow
connection.  Every component is a projection
of~$\Theta_\cA$.
\end{definition}

The datum $\cH(\cA)$ is algorithmically
reconstructed from the boundary OPE of~$\cA$:
$\cA \xmapsto{B} \Theta_\cA \xmapsto{\Res^{\mathrm{coll}}}
r(z) \xmapsto{\mathrm{RTT}} \mathcal{C}$.  Each step uses
only the output of the preceding step; the entire holographic
system is determined by the chiral algebra on the boundary.

% ====================================================================
\subsection{The 4d--3d--2d resolution}
\label{ssec:4d-3d-2d}
% ====================================================================

The three volumes of the programme correspond to three levels
of a dimensional hierarchy:

\begin{center}
\renewcommand{\arraystretch}{1.3}
\begin{tabular}{cllll}
\toprule
\textbf{$d$} & \textbf{Level} & \textbf{Operadic type}
  & \textbf{Key datum} & \textbf{Volume} \\
\midrule
$4$ & CY source & $E_2$
  & CY category $\mathcal{C}$ & III \\
$3$ & Swiss-cheese & $E_1$
  & $R$-matrix $r(z)$ & II \\
$2$ & Modular shadow & $E_\infty$
  & $\kappa(\cA)$ & I \\
\bottomrule
\end{tabular}
\end{center}

\noindent
At each level, information is lost.  From 4d to 3d, one
topological direction is forgotten: $E_2 \to E_1$ (the
braiding becomes ordering).  From 3d to 2d,
$\Sigma_n$-coinvariance projects the $R$-matrix to a scalar:
$r(z) \mapsto \kappa$.  The five main theorems are the
invariants that survive both projections.

% ====================================================================
\subsection{Gauge versus gravity}
\label{ssec:gauge-vs-gravity}
% ====================================================================

The shadow depth classification
(Definition~\ref{def:shadow-depth}) separates the standard
landscape into two physical regimes:

\begin{enumerate}[(a)]
\item \emph{Gauge theories} (classes~$\mathbf{G}$,
  $\mathbf{L}$): finite shadow depth, Swiss-cheese-formal bar
  complex ($m_k^{\mathrm{SC}} = 0$ for $k$ sufficiently large).
  The collision residue has at most a simple pole.
  The bar complex is effectively quadratic or cubic.
  The representation theory is governed by quantum groups with
  rational $R$-matrices.

\item \emph{Gravitational theories} (classes~$\mathbf{C}$,
  $\mathbf{M}$): the shadow tower extends to arity~$4$
  (class~$\mathbf{C}$) or is infinite (class~$\mathbf{M}$).
  The collision residue has higher-order poles.  The bar complex
  is genuinely non-formal: the full $A_\infty$-tower is
  irreducible.  The Koszul duality
  $\Vir_c^! = \Vir_{26-c}$ (or
  $\cW_N^! = \cW_N^{k'}$ at the dual level)
  exchanges matter and ghost sectors.
\end{enumerate}

DS reduction is the functor connecting these regimes: it
transports affine Kac--Moody (class~$\mathbf{L}$, gauge) to
$\cW$-algebras (class~$\mathbf{M}$, gravity) by
raising the pole order through the Sugawara construction.
The physical content: the Sugawara stress tensor converts
gauge currents into a metric, and the bar complex detects
this conversion as the appearance of higher collision poles.

% ====================================================================
\subsection{The critical string at $c = 26$}
\label{ssec:critical-string}
% ====================================================================

The Virasoro Koszul duality $\Vir_c^! = \Vir_{26 - c}$
identifies three distinguished points:

\begin{center}
\renewcommand{\arraystretch}{1.3}
\begin{tabular}{clll}
\toprule
$\boldsymbol{c}$ & \textbf{Property}
  & $\boldsymbol{\kappa}$ & \textbf{Clifford} \\
\midrule
$0$ & Gauge point: $\kappa = 0$, flat bar
  & $0$ & degenerate \\
$13$ & Self-dual point: $\Vir_{13}^! = \Vir_{13}$
  & $13/2$ & non-degenerate \\
$26$ & Critical string: $\kappa^! = 0$
  & $13$ & exterior algebra \\
\bottomrule
\end{tabular}
\end{center}

\noindent
At $c = 0$: $\kappa(\Vir_0) = 0$, the bar complex is flat at
all genera, and the genus tower is trivial.  At $c = 13$:
$\kappa(\Vir_{13}) = 13/2 = \kappa(\Vir_{13}^!)$, the
complementarity pairing is symmetric, and the entire shadow
tower is self-dual.  At $c = 26$: $\kappa(\Vir_0^!) = 0$, the
Koszul dual is flat; the Clifford algebra of the
complementarity pairing degenerates to an exterior algebra,
and the matter-ghost cancellation condition
$\kappa_{\mathrm{eff}}
= \kappa(\mathrm{matter}) + \kappa(\mathrm{ghost}) = 0$ of the
bosonic string is satisfied.  The $26$-dimensional critical
string is the point where the Koszul dual's bar complex
becomes flat: the gravitational content migrates from curvature
to bar cohomology.


% ====================================================================
\section{The Calabi--Yau bridge}
\label{sec:cy-bridge}
% ====================================================================

Volumes~I and~II take a chiral algebra as given.  The source
of chiral algebras is geometry: a Calabi--Yau category
$\mathcal{C}$ of dimension~$d$ produces a chiral algebra
$A_\mathcal{C}$ on a curve~$X$, with the CY trace realized
as the modular characteristic.  This is the content of
Volume~III.

% ====================================================================
\subsection{The functor $\Phi$}
\label{ssec:functor-phi}
% ====================================================================

The construction proceeds in three steps.  Given a CY$_d$
category $\mathcal{C}$:

\begin{enumerate}[(1)]
\item The cyclic $A_\infty$-structure on $\mathcal{C}$
  determines a \emph{Lie conformal algebra} $L_\mathcal{C}$:
  the cyclic bar differential gives the $\lambda$-bracket,
  the cyclic trace gives the invariant bilinear form, and
  the higher $A_\infty$-operations give higher Lie conformal
  products.

\item The \emph{factorization envelope}
  $U_X^{\mathrm{fact}}(L_\mathcal{C})$ produces a
  factorization algebra on $\Ran(X)$, hence a chiral algebra
  $A_\mathcal{C}$ on~$X$.  The PBW theorem for factorization
  envelopes guarantees the expected filtration.

\item The $\mathbb{S}^d$-framing on Hochschild homology
  determines the operadic enhancement.  For $d = 2$: the
  $\mathbb{S}^2$-action gives an $E_2$-structure with braided
  monoidal representation category.  For $d = 3$: holomorphic
  Chern--Simons theory on $\mathcal{C}$ breaks $E_2$ to $E_1$,
  and the braided structure is recovered via the Drinfeld centre
  of the $E_1$-monoidal category.
\end{enumerate}

\begin{theorem}[CY-A$_2$]
\label{thm:cy-a2}
The functor $\Phi$ is constructed for $d = 2$: all three steps
are proved.  The modular characteristic of $A_\mathcal{C}$
equals the Euler characteristic:
$\kappa(A_\mathcal{C}) = \chi^{\mathrm{CY}}(\mathcal{C})$.
\end{theorem}

For $d = 3$, steps~(1) and~(2) are unconditional; step~(3) is
a programme conditional on the chain-level
$\mathbb{S}^3$-framing.  The passage from $d = 2$ to $d = 3$
replaces a geometric braiding (from
$\pi_1(\Conf_2(\mathbb{R}^2)) = \mathbb{Z}$) with an
algebraic one (from the Drinfeld centre of the
$E_1$-monoidal representation category): this is the CY
analogue of the Tannakian reconstruction of a quantum group
from its module category.

% ====================================================================
\subsection{The $E_1$ stabilization}
\label{ssec:e1-stabilization}
% ====================================================================

For $d \ge 3$, the CY chiral algebra is $E_1$, not $E_d$.
The reason is structural: the Schouten--Nijenhuis bracket on
polyvector fields of $\Bbbk^d$ vanishes identically for
$d \ge 3$, so the classical Lie conformal algebra is abelian
and all noncommutative structure arises from quantization.
The CoHA construction produces an ordered (associative) product
from the extension correspondence
$0 \to V' \to V \to V'' \to 0$, which has a preferred
direction: sub before quotient.  This is $E_1$, not $E_2$.

The framing obstruction lies in $\pi_d(BU)$, which is
$\mathbb{Z}$ for even~$d$ and $0$ for $d \equiv 1, 3, 7
\pmod{8}$.  The pattern repeats with period~$8$ by Bott
periodicity.  For $d = 2$: the obstruction is
$c_1 \in \pi_2(BU) = \mathbb{Z}$, the braiding parameter.
For $d = 3$: the obstruction vanishes
($\pi_3(BU) = 0$), so the $\mathbb{S}^3$-framing trivializes
and the $E_1$-structure receives no additional enhancement.
For $d = 4$: $\pi_4(BU) = \mathbb{Z}$ produces a shifted
symplectic level ($p_1$), the CY$_4$ analogue of the
Kac--Moody level.

% ====================================================================
\subsection{The $K3 \times E$ prototype}
\label{ssec:k3-times-e}
% ====================================================================

The product of a K3 surface~$S$ and an elliptic curve~$E$
is the canonical CY$_3$ testing ground: simple enough to
compute, rich enough to obstruct.

The lattice $\Lambda^{3,2}$ has signature $(3,2)$; the
Jacobi form $\phi_{0,1}(\tau, z)$ is the K3 elliptic genus;
the Igusa cusp form $\Delta_5$ is the Borcherds denominator
identity.  The modular characteristic of the chiral de Rham
complex is $\kappa_{\mathrm{ch}}(K3 \times E) = 3
= \dim_\Bbbk$, verified by six independent paths.  The BKM
automorphic weight is a distinct quantity:
$\kappa_{\mathrm{BKM}} = 5$ (the weight of $\Delta_5$).

These are different numbers measuring different things.
A single CY$_3$ admits multiple chiral algebraizations,
each producing a different modular characteristic.  The
$\kappa$-spectrum
\begin{equation}\label{eq:kappa-spectrum}
\operatorname{Spec}_\kappa(K3 \times E)
\;=\; \{2, 3, 5, 24\}
\end{equation}
records four values: $\kappa_{\mathrm{cat}} = 2
= \chi(\mathcal{O}_{K3})$ (the holomorphic Euler
characteristic), $\kappa_{\mathrm{ch}} = 3$ (the complex
dimension), $\kappa_{\mathrm{BKM}} = 5$ (the Borcherds
weight), and $\kappa_{\mathrm{fiber}} = 24$ (the lattice
rank).  Four algebras see four aspects of the same manifold.

Four structural obstructions separate the shadow obstruction
tower from the full automorphic form:
\begin{enumerate}[label=\textup{(O\arabic*)}]
\item \emph{Categorical}: the shadow tower produces Taylor
  coefficients; the denominator identity is an automorphic form
  on $\cH_2$.
\item \emph{$\kappa$-mismatch}:
  $\kappa_{\mathrm{ch}} = 3 \ne 5 = \kappa_{\mathrm{BKM}}$.
\item \emph{Second quantization}: the physical DT partition
  function is the DMVV symmetric product of single-copy data;
  the bar complex is the single copy.
\item \emph{Schottky}: at genus $g \ge 4$, the Torelli map
  $\cM_g \to \mathcal{A}_g$ is no longer dominant, and the
  shadow tower is imprisoned in tautological classes.
\end{enumerate}
Each obstruction points to a research programme.  The
$\kappa$-spectrum is the first invariant that separates CY
chiral algebras with the same underlying manifold.


% ====================================================================
\subsection{$E_1$ at $d = 3$ via holomorphic Chern--Simons}
\label{ssec:e1-holcs}
% ====================================================================

For a CY threefold $X$, holomorphic Chern--Simons theory
on $X$ with gauge group~$G$ produces a chiral algebra
$A_X$ of $E_1$-type.  The spectral parameter of the
$R$-matrix is identified with the holomorphic coordinate
of the 4d theory restricted to a 2d boundary.  The
braided monoidal structure on the representation category is
recovered through the Drinfeld centre:
\begin{equation}\label{eq:drinfeld-centre}
\mathcal{Z}\bigl(\mathrm{Rep}^{E_1}(A_X)\bigr)
\;\simeq\;
\mathrm{Rep}^{E_2}(A_X)\,,
\end{equation}
where the right-hand side is the braided monoidal category of
$E_2$-modules.  For toric CY$_3$, the $E_1$-algebra is the
critical CoHA $\mathrm{CoHA}(Q_X)$: the positive half of an
affine super Yangian $Y(\widehat{\fg}_{Q_X})$.  For
$\Bbbk^3$: the Jordan quiver gives
$Y(\widehat{\mathfrak{gl}}_1) \simeq \cW_{1+\infty}$.

The dimensional reduction principle underlies the full trilogy:
$4d \to 3d$ forgets the braiding ($E_2 \to E_1$); $3d \to 2d$
forgets the ordering ($E_1 \to E_\infty$ by
$\Sigma_n$-coinvariance).  The five main theorems live at the
$2d$ level; the seven faces of the collision residue live at
the $3d$ level; the CY source lives at the $4d$ level.  The
modular Koszul engine converts enumerative geometry
(Donaldson--Thomas invariants, root multiplicities, denominator
identities) into homotopy algebra (bar complexes, MC elements,
genus expansions).  One Maurer--Cartan element
$\Theta_{A_\mathcal{C}}$ absorbs the entire automorphic datum.


% ====================================================================
\section{The standard landscape}
\label{sec:landscape}
% ====================================================================

The bar construction, the genus tower, the Swiss-cheese
structure, and the collision residue are not abstractions
awaiting examples: they were \emph{found} by computing
examples.  The standard landscape is the testing ground
where the theory proves itself.

% ====================================================================
\subsection{Census}
\label{ssec:census}
% ====================================================================

Seven families span the landscape.  Each is chirally
Koszul.  The invariants computed in the monograph are
recorded in Table~\ref{tab:census}.

\begin{table}[ht]
\centering
\caption{The standard landscape: seven families and their
  bar-complex invariants}
\label{tab:census}
\renewcommand{\arraystretch}{1.7}
{\small
\begin{tabular}{lcccccl}
\toprule
\textbf{Algebra $\cA$}
  & $\boldsymbol{c(\cA)}$
  & $\boldsymbol{\kappa(\cA)}$
  & $\boldsymbol{r_{\max}}$
  & \textbf{Class}
  & $\boldsymbol{r(z)}$
  & $\boldsymbol{\cA^!}$ \\
\midrule
$\cH_k$
  & $1$
  & $k$
  & $2$ & $\mathbf{G}$
  & $k/z$
  & $\Sym^{\mathrm{ch}}(V^*)$ \\[2pt]
$\widehat{\fg}_k$
  & $\dfrac{kd}{k+h^\vee}$
  & $\dfrac{d(k+h^\vee)}{2h^\vee}$
  & $3$ & $\mathbf{L}$
  & $k\,\Omega/z$
  & $\widehat{\fg}_{-k-2h^\vee}$ \\[6pt]
$\beta\gamma_\lambda$
  & $2(6\lambda^2{-}6\lambda{+}1)$
  & $c/2$
  & $4$ & $\mathbf{C}$
  & const.
  & $bc_\lambda$ \\[2pt]
$\Vir_c$
  & $c$
  & $c/2$
  & $\infty$ & $\mathbf{M}$
  & $\frac{c/2}{z^3} + \frac{2T}{z}$
  & $\Vir_{26-c}$ \\[4pt]
$\cW^k_3(\mathfrak{sl}_3)$
  & $c(k)$
  & $5c/6$
  & $\infty$ & $\mathbf{M}$
  & $r_{\cW_3}(z)$
  & $\cW^{k'}_3$ \\[2pt]
$\cW^k_N(\mathfrak{sl}_N)$
  & $c(k)$
  & $c \cdot (H_N{-}1)$
  & $\infty$ & $\mathbf{M}$
  & $r_{\cW_N}(z)$
  & $\cW^{k'}_N$ \\[2pt]
$V_\Lambda$\,(lattice)
  & $\mathrm{rank}(\Lambda)$
  & $\mathrm{rank}(\Lambda)$
  & $2$ & $\mathbf{G}$
  & $\mathrm{rank}/z$
  & $V_{\Lambda^*}$ \\
\bottomrule
\end{tabular}
}
\end{table}

\noindent
Here $d = \dim\fg$, $h^\vee$ is the dual Coxeter number,
$t = k + h^\vee$, $k' = -k - 2h^\vee$,
$H_N = \sum_{j=1}^{N} 1/j$ is the $N$-th harmonic number,
and $\Omega = \sum_a J^a \otimes J_a$ is the Casimir tensor.

The table has two readings.  Read horizontally, each row is
a complete portrait: central charge, modular characteristic,
shadow depth, collision residue, Koszul dual.  Read
vertically, the columns expose the structural regularities
that no single example can reveal.  The modular characteristic
$\kappa$ is always a rational function of the level;
the complementarity sum $\kappa + \kappa^!$ vanishes for
free-field and Kac--Moody algebras, equals $13$ for
Virasoro, and is $(c + c')(H_N - 1)$ for
$\cW_N$.  The collision residue $r(z)$ is scalar for
classes~$\mathbf{G}$ and~$\mathbf{C}$, matrix-valued for
classes~$\mathbf{L}$ and~$\mathbf{M}$.

% ====================================================================
\subsection{The DS reduction chain}
\label{ssec:ds-chain}
% ====================================================================

Drinfeld--Sokolov reduction links the families.  For
$\fg = \mathfrak{sl}_N$, the principal DS reduction
\begin{equation}\label{eq:ds-chain}
V_k(\mathfrak{sl}_N)
\;\xrightarrow{\;\mathrm{DS}\;}
\;\cW^k_N(\mathfrak{sl}_N)
\end{equation}
transforms the shadow obstruction tower in a controlled
way:
\begin{enumerate}[(i)]
\item \emph{Depth increase.}
  The affine algebra $V_k(\mathfrak{sl}_N)$ is
  class~$\mathbf{L}$ ($r_{\max} = 3$); the
  $\cW$-algebra $\cW^k_N$ is class~$\mathbf{M}$
  ($r_{\max} = \infty$).  DS reduction creates the
  infinite shadow tower from finite data.

\item \emph{Quartic creation from zero.}
  The quartic shadow $S_4$ vanishes for all affine
  algebras (the Jacobi identity kills the
  obstruction); after DS reduction,
  $S_4(\cW^k_N) \ne 0$ for generic~$k$.
  For Virasoro:
  $S_4(\Vir_c) = 10/[c(5c+22)]$.

\item \emph{Central charge is additive; $\kappa$ is not.}
  The BRST ghost system contributes
  $c_{\mathrm{ghost}} = N(N{-}1)$ to the central
  charge but a $k$-dependent correction to~$\kappa$.
  The commutation failure of DS reduction with the
  shadow map is the source of all non-formality in
  the $\cW$-algebra landscape.
\end{enumerate}

The chain $\cH \to \widehat{\fg}_k \to \cW^k_N$ is a
cascade of increasing shadow complexity:
$\mathbf{G} \to \mathbf{L} \to \mathbf{M}$, with each
reduction introducing new structure that the previous
level could not see.


% ====================================================================
\subsection{The multi-weight correction}
\label{ssec:multi-weight-landscape}
% ====================================================================

For algebras with generators of distinct conformal
weights, the genus expansion
receives a correction.  The $\cW_3$ algebra has generators
$T$ (weight~$2$) and $W$ (weight~$3$); at genus~$2$ the
cross-channel graph sum contributes
\begin{equation}\label{eq:delta-f2-w3}
\delta F_2(\cW_3)
\;=\;
\frac{c + 204}{16c}
\;>\; 0
\qquad
\text{for all } c > 0\,.
\end{equation}
The correction arises from boundary graphs
of~$\overline{\cM}_{2,0}$ whose edges carry mixed
propagator channels (weight-$2$ and weight-$3$ generators
sewed through different Hodge data at the vertex level).
The propagator itself is always weight-$1$ at the edge
level ($d\log E$ has conformal weight~$1$); the mixing
enters through the vertex structure constants.

For uniform-weight algebras (all generators of the same
conformal weight), $\delta F_g^{\mathrm{cross}} = 0$
at all genera and the scalar universality
$F_g = \kappa \cdot \lambda_g^{\mathrm{FP}}$ holds
without correction.  For multi-weight algebras, the
full genus expansion is
$F_g = \kappa \cdot \lambda_g^{\mathrm{FP}}
+ \delta F_g^{\mathrm{cross}}$,
with
$\delta F_g^{\mathrm{cross}}$ computable from the
mixed-channel graph sum.



% ====================================================================
\section{The MC programme}
\label{sec:mc-programme}
% ====================================================================

The MC programme converts the geometric structure of the
bar complex into five concrete mathematical problems, each
asking whether a structural property of the universal MC
element~$\Theta_\cA$ persists under a specific operation:
PBW concentration, bar-intrinsic existence, categorical
generation, inverse-limit completion, and analytic sewing.
MC1 through MC4 are proved; MC5 is partially proved, with
the analytic HS-sewing package established at all genera
and the genuswise BV/BRST/bar identification conjectural.

% ====================================================================
\subsection{MC1: PBW concentration}
\label{ssec:mc1}
% ====================================================================

The PBW spectral sequence of~$B(\cA)$ collapses at~$E_2$
for every freely strongly generated chiral algebra.  Free
strong generation implies the PBW filtration on the bar
complex is exhaustive, and the associated graded is a
cofree coalgebra whose cohomology concentrates in bar
degree~$1$.  This is chirally Koszul: the bar cohomology
$H^*(B(\cA))$ is the Koszul dual $\cA^!$ (as a graded
coalgebra).

The proof applies to the universal $\cW$-algebra
$\cW^k(\fg)$ at all levels~$k$ (Feigin--Frenkel free
generation~\cite{FF92,FBZ04}).  The simple quotient $L_k(\fg)$ at admissible
levels has Koszulness proved for $\fg = \mathfrak{sl}_2$
at all admissible levels; higher rank is open.

% ====================================================================
\subsection{MC2: bar-intrinsic existence}
\label{ssec:mc2}
% ====================================================================

The universal MC element $\Theta_\cA = D_\cA - d_0$
exists for all modular Koszul chiral algebras, with no
restriction to simple Lie symmetry.  The proof is a
tautology at the right level of abstraction: $D_\cA^2 = 0$
implies $\Theta_\cA$ satisfies the MC
equation~(\ref{eq:mc-equation}).  The algebraic-family
rigidity theorem extends the scalar identification
$\Theta_\cA^{\min} = \eta \otimes \Gamma_\cA$ to all
algebraic families with rational OPE coefficients, covering
the entire standard Lie-theoretic landscape at all
non-critical levels.

% ====================================================================
\subsection{MC3: categorical generation}
\label{ssec:mc3}
% ====================================================================

For every simple Lie algebra~$\fg$, the Drinfeld--Kazhdan
category $\mathrm{DK}(\fg)$ is thickly generated by
evaluation modules.  The proof combines the chromatic
filtration on $\mathrm{Rep}(Y(\fg))$ with the categorical
Clebsch--Gordan closure via multiplicity-free
$\ell$-weights (Chari--Moura~\cite{CM08}), the Francis--Gaitsgory
pro-nilpotent completion~\cite{FG12}, and Drinfeld--Kazhdan
compactness~\cite{DK84}.  The type~A proof uses the minuscule
hypothesis; the all-types generalization replaces it with
the strictly weaker multiplicity-free $\ell$-weight
property, which holds for all simple types.

% ====================================================================
\subsection{MC4: inverse-limit completion}
\label{ssec:mc4}
% ====================================================================

The bar-cobar adjunction extends to inverse limits via the
strong completion-tower theorem: the strong filtration axiom
$\mu_r(F^{i_1}, \ldots, F^{i_r})
\subset F^{i_1 + \cdots + i_r}$
forces an arity cutoff that makes continuity and the
Mittag-Leffler condition automatic.  The completion closure
$\mathrm{CompCl}(F_{\mathrm{ft}})$ carries a quasi-inverse
bar-cobar homotopy equivalence, stable under MC twisting.

The programme splits: MC4$^+$ (positive towers:
$\cW_{1+\infty}$, affine Yangians) is solved by weight
stabilization.  MC4$^0$ (resonant towers: Virasoro,
non-quadratic $\cW_N$) is reduced to a finite resonance
problem by the resonance-filtered bar-cobar theorem.

% ====================================================================
\subsection{MC5: analytic sewing}
\label{ssec:mc5}
% ====================================================================

The sewing envelope $\cA^{\mathrm{sew}}$, the Hausdorff
completion of~$\cA$ for the all-amplitude seminorm
topology, carries the analytic data needed for genus-$g$
partition functions.  The HS-sewing criterion provides a
sufficient condition: polynomial OPE growth and
subexponential sector growth imply convergence at all
genera.  Every standard family satisfies this condition.

For Heisenberg, the sewing theorem is explicit: the
one-particle Bergman reduction writes the genus-$g$
partition function as a Fredholm determinant.  For general
modular Koszul algebras, the HS-sewing criterion gives
convergence without an explicit closed form.  The BV/BRST
identification of sewing amplitudes with bar complex
data at higher genus remains conjectural.


% ====================================================================
\section{Arithmetic}
\label{sec:arithmetic}
% ====================================================================

The shadow obstruction tower produces arithmetic invariants
from vertex-algebraic data.  These invariants connect the
bar complex to number theory through three structures: a
Dirichlet series, a categorical zeta function, and a depth
decomposition.

% ====================================================================
\subsection{The shadow Eisenstein theorem}
\label{ssec:eisenstein}
% ====================================================================

The genus-$1$ amplitude of a modular Koszul chiral
algebra~$\cA$ is the quasi-modular Eisenstein series
$\kappa(\cA) \cdot E_2^*(\tau)$, where
$E_2^*(\tau) = E_2(\tau) - 3/(\pi \operatorname{Im}\tau)$
is the completed weight-$2$ Eisenstein series.  The
Fourier coefficients $\sigma_1(n)$ of $E_2^*$ produce
a Dirichlet series:
\begin{equation}\label{eq:dirichlet-series}
D_2(\cA, s)
\;:=\;
-24\kappa(\cA) \cdot
\zeta(s) \cdot \zeta(s - 1)\,,
\end{equation}
an Eisenstein $L$-function with poles at $s = 1$ (from
$\zeta(s)$) and $s = 2$ (from $\zeta(s-1)$).  The identity
follows from the $\sigma_1$ Dirichlet convolution applied to
the Fourier coefficients of $\kappa \cdot E_2^*(\tau)$.

The genus-$1$ propagator is $E_2^*$, not a holomorphic
modular form: $E_2^*$ is quasi-modular of weight~$2$ and
depth~$1$, transforming with an additive anomaly
$3\tau/(\pi i)$ under $\mathrm{SL}(2,\mathbb{Z})$.
Products of $E_2^*$ live in the ring of quasi-modular
forms $\mathbb{C}[E_2^*, E_4, E_6]$, not in the ring of
holomorphic modular forms $\mathbb{C}[E_4, E_6]$.  This
distinction is load-bearing: holomorphic dimension formulas
(e.g., $\dim S_k = 0$ for $k < 12$) do not apply to
quasi-modular objects.


% ====================================================================
\subsection{The categorical zeta function}
\label{ssec:categorical-zeta}
% ====================================================================

The dimension-counting zeta function of
$\mathfrak{sl}_2$-representations is
\begin{equation}\label{eq:cat-zeta}
\zeta^{\mathrm{DK}}_{\mathfrak{sl}_2}(s)
\;:=\;
\sum_{n \ge 1} (\dim V_n)^{-s}
\;=\;
\sum_{n \ge 1}(n+1)^{-s}
\;=\;
\zeta(s) - 1\,.
\end{equation}
The Riemann zeta function (shifted by~$1$) counts
$\mathfrak{sl}_2$-representations weighted by inverse
dimension.  The link to the shadow--DK bridge is that
the categorical zeta encodes the multiplicity structure
of the evaluation modules whose thick generation proves
MC3.

For $\mathfrak{sl}_N$ with $N \ge 3$,
$\zeta^{\mathrm{DK}}_{\mathfrak{sl}_N}(s) :=
\sum_\lambda (\dim V_\lambda)^{-s}$ is not
multiplicative: the Weyl dimension formula is a product
of linear forms, but the sum over the dominant chamber
has no Euler product.  The $N = 2$ coincidence with
$\zeta(s)$ is a low-rank accident.

% ====================================================================
\subsection{Depth decomposition}
\label{ssec:depth}
% ====================================================================

The total shadow depth decomposes:
\begin{equation}\label{eq:depth-decomposition}
d(\cA)
\;=\;
1 + d_{\mathrm{arith}}(\cA) + d_{\mathrm{alg}}(\cA)\,,
\end{equation}
where $d_{\mathrm{arith}}$ counts arithmetic contributions
(cusp forms in the automorphic decomposition of the shadow
generating function) and $d_{\mathrm{alg}}$ counts
algebraic contributions (non-formality of the $L_\infty$
structure).  The algebraic depth
$d_{\mathrm{alg}} \in \{0, 1, 2, \infty\}$ recovers
the G/L/C/M classification: $d_{\mathrm{alg}} = 0$
for class~$\mathbf{G}$, $d_{\mathrm{alg}} = 1$ for
class~$\mathbf{L}$, $d_{\mathrm{alg}} = 2$ for
class~$\mathbf{C}$,
$d_{\mathrm{alg}} = \infty$ for class~$\mathbf{M}$.
Shadow depths $\ge 5$ are purely arithmetic: they arise
from cusp forms, beginning with the weight-$12$ cusp
form $\Delta_{12}$ at the genus-$2$ level.

The decomposition separates the algebraic content
(governed by the MC equation on the cyclic deformation
complex) from the arithmetic content (governed by
automorphic forms on moduli of curves).  The two are
independent: a class-$\mathbf{G}$ algebra (algebraically
trivial) can have nontrivial arithmetic depth, and a
class-$\mathbf{M}$ algebra (algebraically infinite)
can have vanishing arithmetic depth at low genus.


% ====================================================================
\section{Frontier directions}
\label{sec:frontier}
% ====================================================================

Three leaps organize the frontier.  The first extends the
modular Koszul programme from chiral algebras to
holographic systems.  The second extends it from formal
algebraic data to analytic sewing amplitudes.  The third
extends it from $2$-dimensional geometry to
$3$-dimensional topology.

% ====================================================================
\subsection{First leap: holographic modular Koszul data}
\label{ssec:holographic-frontier}
% ====================================================================

Every holomorphic-topological system~$T$ in the sense of
Costello--Li~\cite{CL16} should be controlled by a \emph{holographic
modular Koszul datum}
\begin{equation}\label{eq:holographic-datum-frontier}
\cH(T)
\;=\;
(\cA,\, \cA^!,\, \mathfrak{C},\, r(z),\,
\Theta_\cA,\, \nabla^{\mathrm{hol}})\,,
\end{equation}
consisting of the boundary algebra, its Koszul dual, the
bulk algebra (chiral derived centre), the collision
residue, the universal MC element, and the holographic
shadow connection.  The collision residue is the genus-$0$
binary shadow of~$\Theta_\cA$:
$r(z) = \Res^{\mathrm{coll}}_{0,2}(\Theta_\cA)$.

This is the Platonic holographic programme.  Its five
theorem targets are: boundary-defect realization (the
boundary algebra is the algebra of boundary conditions
of the $3$d theory), Yangian-shadow identification (the
Yangian controls line operators), sphere reconstruction
(genus-$0$ amplitudes recover the boundary OPE), quartic
resonance obstruction (the quartic shadow controls the
deformation quantization), and singular-fibre descent (the
critical-level limit produces the Feigin--Frenkel centre).

Two other directions in the first leap: the
factorization-envelope technology of Nishinaka~\cite{Nis25} and Vicedo~\cite{Vic25},
which gives a functorial pipeline from Lie conformal
algebras to vertex algebras; and the non-principal
$\cW$-algebra corridor, where hook-type nilpotent orbits
in type~A provide the first proved non-principal Koszul
dualities.


% ====================================================================
\subsection{Second leap: analytic sewing}
\label{ssec:sewing-frontier}
% ====================================================================

The algebraic bar construction produces formal power
series in the sewing parameters.  The analytic sewing
programme asks: when do these series converge?

The answer is the sewing envelope $\cA^{\mathrm{sew}}$,
the Hausdorff completion for the all-amplitude seminorm
topology.  The HS-sewing condition (polynomial OPE growth
and subexponential sector growth) is satisfied by the
entire standard landscape.  The Heisenberg sewing theorem
gives an explicit Fredholm determinant.  The lattice
sewing envelope is constructed; the analytic realization
criterion (a sufficient condition for passage from
formal to analytic) is conjectural.

At genus $g \ge 1$ the curvature~\eqref{eq:curvature}
forces passage from ordinary to coderived categories:
the bar complex is no longer flat, and the correct
homological framework is the coderived category of
curved coalgebras.  The analytic programme at higher genus
requires the coderived analytic shadow
$Q_g^{\mathrm{an}}(\cA)$.


% ====================================================================
\subsection{Third leap: three-dimensional topology}
\label{ssec:3d}
% ====================================================================

The Swiss-cheese operad $\mathrm{SC}^{\mathrm{ch,top}}$
packages both the chiral (holomorphic, $\mathbb{C}$)
and topological ($\mathbb{R}$) directions of a
$3$-dimensional holomorphic-topological theory.  The
bar complex factorizes along both: the differential
encodes the $\mathbb{C}$-direction factorization, the
coproduct encodes the $\mathbb{R}$-direction
factorization.  The ordered bar
$\barB^{\mathrm{ord}}(\cA)$ is simultaneously an
$E_1$-coalgebra (from the linear order on~$\mathbb{R}$)
and an $E_\infty$-coalgebra (from the factorization
structure on~$\mathbb{C}$).

The third leap is the passage from the algebraic
Swiss-cheese structure (proved) to the geometric
realization on $3$-manifolds of the form
$\Sigma_g \times [0,1]$.  The $\mathbb{C}$-direction
carries the chiral algebra; the $\mathbb{R}$-direction
carries the Yangian / quantum group.  The annulus trace
$\Delta_{\mathrm{ns}}(\mathrm{Tr}_\cA) =
\kappa \cdot \lambda_1$
provides the first open-to-closed map at genus~$1$;
the full genus tower progressively restores
bidirectionality between open and closed sectors.

The deepest open problem in this direction is CY-A at
$d = 3$ (Section~\ref{sec:open-problem}).


% ====================================================================
\section{The single open problem}
\label{sec:open-problem}
% ====================================================================

Everything in Sections~\ref{sec:bar}--\ref{sec:frontier}
rests on two-dimensional geometry: curves, their moduli
spaces, and the factorization structure of colliding points
on a curve.  The three-dimensional extension is conditional
on a single unresolved input.

\begin{conjecture}[CY-A at $d = 3$]\label{conj:cy-a}
The Calabi--Yau/$A_\infty$ correspondence extends from
$d = 2$ (proved) to $d = 3$: the chain-level
$S^3$-framing of the topological direction in a
$3$-dimensional holomorphic-topological theory produces
an $A_\infty$-structure on the boundary algebra that is
compatible with the chiral bar complex.
\end{conjecture}

At $d = 2$, the CY-A correspondence is a theorem: the
Fukaya category of a symplectic surface carries an
$A_\infty$-structure whose bar construction is the wrapped
Floer complex, and the $A_\infty$-formality on the Koszul
locus is the $E_2$-collapse of the bar spectral sequence.
At $d = 3$, the analogue requires the chain-level framing
of the $S^3$-fibre in the conifold geometry to be
compatible with the factorization structure on the
$\mathbb{C}$-direction.  This is a statement about the
interaction of holomorphic and topological data at the
chain level, not merely at the level of cohomology.

The difficulty is not conceptual but technical: the
$d = 2$ proof uses the Riemann surface structure in an
essential way (the Abel--Jacobi map, the period matrix,
the prime form), and the $d = 3$ analogue must replace
surface geometry with $3$-manifold geometry while
preserving the $A_\infty$-coherences.

No other open conjecture in the programme has unresolved
algebraic prerequisites.  The five MC problems are solved;
the Koszulness characterization is complete to within
the D-module purity converse; the shadow obstruction tower
is constructed at all arities; the arithmetic invariants
are computed; the standard landscape is fully mapped.
Six new results extend the proved core: the universal
genus-$g$ shadow identity $\alpha_g = \kappa \cdot
\lambda_g^{\mathrm{FP}}$ at all genera (uniform-weight),
the soft graviton theorem $a_3 = 2$ fixing the cubic
shadow coefficient, and free-field exactness of the bar
complex (the bar spectral sequence collapses at $E_1$
for all free-field algebras).  The full programme now
comprises 4,542 pages across three volumes, verified by
120K+ independent computational tests.
CY-A at $d = 3$ is the single point where the programme
reaches beyond what it can currently prove, and it is the
natural boundary: the passage from the algebraic curve
to the three-manifold.


% ====================================================================
\section*{References}
\label{sec:references}
% ====================================================================

\begin{thebibliography}{99}

\bibitem{Arnold69}
V.\,I.~Arnold,
The cohomology ring of the group of dyed braids,
\textit{Mat.\ Zametki}~\textbf{5} (1969), 227--231.

\bibitem{BD04}
A.~Beilinson and V.~Drinfeld,
\textit{Chiral Algebras},
AMS Colloquium Publications, vol.~51,
American Mathematical Society, Providence, RI, 2004.

\bibitem{BGS96}
A.~Beilinson, V.~Ginzburg, and W.~Soergel,
Koszul duality patterns in representation theory,
\textit{J.~Amer.\ Math.\ Soc.}~\textbf{9} (1996),
473--527.

\bibitem{CG17}
K.~Costello and O.~Gwilliam,
\textit{Factorization Algebras in Quantum Field Theory},
vols.~1--2, Cambridge University Press, 2017/2021.

\bibitem{CL16}
K.~Costello and S.~Li,
Twisted supergravity and its quantization,
preprint, \texttt{arXiv:1606.00365}, 2016.

\bibitem{CM08}
V.~Chari and A.~Moura,
Characters and blocks for finite-dimensional
representations of quantum affine algebras,
\textit{Int.\ Math.\ Res.\ Not.}\ (2004), no.~5,
257--298.

\bibitem{DK84}
V.\,G.~Drinfeld and V.\,G.~Kazhdan,
Algebraic geometry of the space of quasi-maps
and the structure of representations of quantum groups,
unpublished notes, 1984.

\bibitem{Dri85}
V.\,G.~Drinfeld,
Hopf algebras and the quantum Yang--Baxter equation,
\textit{Dokl.\ Akad.\ Nauk SSSR}~\textbf{283} (1985),
1060--1064.

\bibitem{FBZ04}
E.~Frenkel and D.~Ben-Zvi,
\textit{Vertex Algebras and Algebraic Curves},
2nd ed., Mathematical Surveys and Monographs, vol.~88,
American Mathematical Society, Providence, RI, 2004.

\bibitem{FF92}
B.~Feigin and E.~Frenkel,
Affine Kac--Moody algebras at the critical level and
Gel'fand--Diki\u{\i} algebras,
\textit{Int.\ J.\ Mod.\ Phys.\ A}~\textbf{7}
(suppl.\ 1A) (1992), 197--215.

\bibitem{FG12}
J.~Francis and D.~Gaitsgory,
Chiral Koszul duality,
\textit{Selecta Math.\ (N.S.)}~\textbf{18} (2012),
27--87.

\bibitem{FM94}
W.~Fulton and R.~MacPherson,
A compactification of configuration spaces,
\textit{Ann.\ of Math.\ (2)}~\textbf{139} (1994),
183--225.

\bibitem{FP00}
C.~Faber and R.~Pandharipande,
Hodge integrals and moduli spaces of curves,
\textit{Enseign.\ Math.\ (2)}~\textbf{48} (2002),
173--195.

\bibitem{KRW03}
V.~Kac, S.-S.~Roan, and M.~Wakimoto,
Quantum reduction for affine superalgebras,
\textit{Comm.\ Math.\ Phys.}~\textbf{241} (2003),
307--342.

\bibitem{LV12}
J.-L.~Loday and B.~Vallette,
\textit{Algebraic Operads},
Grundlehren der Mathematischen Wissenschaften,
vol.~346, Springer, Berlin, 2012.

\bibitem{Mok25}
C.-P.~Mok,
\emph{Logarithmic Fulton--MacPherson configuration spaces},
arXiv:2503.17563, 2025.

\bibitem{MS24}
N.~Markarian and H.~L.~Tanaka,
Homotopy chiral algebras,
preprint, 2024.

\bibitem{Nis25}
D.~Nishinaka,
Factorization envelopes of Lie conformal algebras,
preprint, 2025.

\bibitem{Priddy70}
S.~Priddy,
Koszul resolutions,
\textit{Trans.\ Amer.\ Math.\ Soc.}~\textbf{152}
(1970), 39--60.

\bibitem{PTVV13}
T.~Pantev, B.~To\"en, M.~Vaqui\'e, and G.~Vezzosi,
Shifted symplectic structures,
\textit{Publ.\ Math.\ Inst.\ Hautes \'Etudes Sci.}
\textbf{117} (2013), 271--328.

\bibitem{RNW19}
D.~Robert-Nicoud and F.~Wierstra,
Homotopy morphisms between convolution homotopy
Lie algebras,
\textit{J.\ Noncommut.\ Geom.}~\textbf{13} (2019),
1463--1520.

\bibitem{Val16}
B.~Vallette,
Homotopy theory of homotopy algebras,
\textit{Ann.\ Inst.\ Fourier}~\textbf{70} (2020),
683--738.

\bibitem{Vic25}
B.~Vicedo,
Non-chiral factorization algebras and vertex algebras,
preprint, 2025.

\end{thebibliography}

\bigskip
\noindent\hrulefill

\medskip
\noindent
Modular Koszul duality for chiral algebras is $E_1$--$E_1$
operadic Koszul duality transported into the homotopical
modular chiral realm on algebraic curves.

\end{document}
